\documentclass[a4paper,11pt]{article}
% Shared preamble for main.tex and the standalone chapter wrappers.
% Hyperref is loaded last. A wrapper defines \chapterbuild before
% % Shared preamble for main.tex and the standalone chapter wrappers.
% Hyperref is loaded last. A wrapper defines \chapterbuild before
% % Shared preamble for main.tex and the standalone chapter wrappers.
% Hyperref is loaded last. A wrapper defines \chapterbuild before
% \input{preamble}; xr-hyper then runs before hyperref and imports
% labels from the full document (default: build/main.aux).
\usepackage[T2A]{fontenc}
\usepackage[utf8]{inputenc}
\usepackage[russian]{babel}
\usepackage{amsmath}
\usepackage{amssymb}

\numberwithin{equation}{section}

\newcommand{\rf}{\mathrm{rf}}
\newcommand{\eq}{\mathrm{eq}}

\ifdefined\chapterbuild
  \usepackage{xr-hyper}
  \providecommand{\MainDoc}{build/main}
  \makeatletter
  \newcommand{\EE@useexternal}[1]{\externaldocument[][nocite]{#1}}
  \expandafter\EE@useexternal\expandafter{\MainDoc}
  \makeatother
\fi
\usepackage[hidelinks]{hyperref}
\ifdefined\chapterbuild
  \makeatletter
  % Aux is read after this hook. Own labels replace the imported ones
  % silently; labels from other chapters stay as in the full document.
  \let\EE@orig@newl@bel\@newl@bel
  \def\EE@silent@newl@bel#1#2#3{{%
    \global\@namedef{#1@#2}{#3}%
  }}
  \AddToHook{begindocument/before}{\let\@newl@bel\EE@silent@newl@bel}
  \AddToHook{begindocument}{\let\@newl@bel\EE@orig@newl@bel}
  \makeatother
\fi
; xr-hyper then runs before hyperref and imports
% labels from the full document (default: build/main.aux).
\usepackage[T2A]{fontenc}
\usepackage[utf8]{inputenc}
\usepackage[russian]{babel}
\usepackage{amsmath}
\usepackage{amssymb}

\numberwithin{equation}{section}

\newcommand{\rf}{\mathrm{rf}}
\newcommand{\eq}{\mathrm{eq}}

\ifdefined\chapterbuild
  \usepackage{xr-hyper}
  \providecommand{\MainDoc}{build/main}
  \makeatletter
  \newcommand{\EE@useexternal}[1]{\externaldocument[][nocite]{#1}}
  \expandafter\EE@useexternal\expandafter{\MainDoc}
  \makeatother
\fi
\usepackage[hidelinks]{hyperref}
\ifdefined\chapterbuild
  \makeatletter
  % Aux is read after this hook. Own labels replace the imported ones
  % silently; labels from other chapters stay as in the full document.
  \let\EE@orig@newl@bel\@newl@bel
  \def\EE@silent@newl@bel#1#2#3{{%
    \global\@namedef{#1@#2}{#3}%
  }}
  \AddToHook{begindocument/before}{\let\@newl@bel\EE@silent@newl@bel}
  \AddToHook{begindocument}{\let\@newl@bel\EE@orig@newl@bel}
  \makeatother
\fi
; xr-hyper then runs before hyperref and imports
% labels from the full document (default: build/main.aux).
\usepackage[T2A]{fontenc}
\usepackage[utf8]{inputenc}
\usepackage[russian]{babel}
\usepackage{amsmath}
\usepackage{amssymb}

\numberwithin{equation}{section}

\newcommand{\rf}{\mathrm{rf}}
\newcommand{\eq}{\mathrm{eq}}

\ifdefined\chapterbuild
  \usepackage{xr-hyper}
  \providecommand{\MainDoc}{build/main}
  \makeatletter
  \newcommand{\EE@useexternal}[1]{\externaldocument[][nocite]{#1}}
  \expandafter\EE@useexternal\expandafter{\MainDoc}
  \makeatother
\fi
\usepackage[hidelinks]{hyperref}
\ifdefined\chapterbuild
  \makeatletter
  % Aux is read after this hook. Own labels replace the imported ones
  % silently; labels from other chapters stay as in the full document.
  \let\EE@orig@newl@bel\@newl@bel
  \def\EE@silent@newl@bel#1#2#3{{%
    \global\@namedef{#1@#2}{#3}%
  }}
  \AddToHook{begindocument/before}{\let\@newl@bel\EE@silent@newl@bel}
  \AddToHook{begindocument}{\let\@newl@bel\EE@orig@newl@bel}
  \makeatother
\fi


\title{Равновесный сдвиг импульса и энергии в накопительном кольце}
\author{}
\date{}

\begin{document}
\maketitle

\begin{abstract}
В заметке решены продольные уравнения движения частицы и выполнено усреднение по периоду синхротронных колебаний. Ключевую роль играет запаздывание частицы на резонаторе относительно синхронной. Сначала рассмотрено магнитное кольцо без учёта бетатронных колебаний. Затем учтена добавка к запаздыванию от бетатронного движения и следующий член разложения ВЧ-напряжения по фазе, который даёт небольшую поправку. Для электростатического кольца предложена модель потенциала электростатического дефлектора, и продольные уравнения решены для тех же случаев. Решения для магнитного и электростатического колец сопоставлены.
\end{abstract}

\tableofcontents

\clearpage
\section{Магнитное кольцо: исходная система}
\label{sec:mag-base}

\subsection{Обозначения}

Индекс $s$ отмечает синхронную частицу: импульс $p_s$, скорость $v_s=\beta_s c$, лоренц-фактор $\gamma_s$, энергия $E_s=\gamma_s mc^2$, длина орбиты $C_s$, период обращения $T_s=C_s/v_s$.

Отклонения импульса и энергии частицы:
\begin{equation}
\label{eq:delta-def}
\delta=\frac{p-p_s}{p_s},
\qquad
w=\frac{E-E_s}{\beta_s^2E_s}.
\end{equation}
Из $\mathrm{d}E/\mathrm{d}p=v$ и $\mathrm{d}^2E/\mathrm{d}p^2=c^2/(\gamma^2E)$
\begin{equation}
\label{eq:w-series}
w=\delta+\frac{\delta^2}{2\gamma_s^2}+O(\delta^3),
\qquad
\delta=w-\frac{w^2}{2\gamma_s^2}+O(w^3).
\end{equation}

Фаза $\varphi$ --- фаза ВЧ-поля в момент прохождения частицей зазора, отсчитанная от синхронной частицы; она растёт, когда частица опаздывает. Кратность $h$, частота $\omega_\rf=2\pi h/T_s$, напряжение зазора $V\sin(\varphi_s+\varphi)$, заряд $q$. Номер оборота $k$, время $t=kT_s$. Относительное запаздывание за оборот
\begin{equation}
\label{eq:tau-def}
\tau=\frac{T-T_s}{T_s},
\end{equation}
где $T$ --- период обращения частицы. Рассматривается одна частица, её траектория лежит внутри сепаратрисы.

\subsection{Запаздывание за оборот}

В магнитном поле скорость вдоль траектории постоянна, поэтому $T=C/v$ и точно
\begin{equation}
\label{eq:T-ratio}
\frac{T}{T_s}=\frac{C}{C_s}\,\frac{v_s}{v}.
\end{equation}
Длина замкнутой орбиты записывается рядом
\begin{equation}
\label{eq:C}
\frac{C}{C_s}=1+\alpha_0\delta+\alpha_1\delta^2+O(\delta^3).
\end{equation}
Из $\beta\gamma=\beta_s\gamma_s(1+\delta)$
\begin{equation}
\label{eq:vsv}
\frac{v_s}{v}=\frac{\gamma}{\gamma_s(1+\delta)}
=1-\frac{\delta}{\gamma_s^2}+\Bigl(1+\frac{\beta_s^2}{2}\Bigr)\frac{\delta^2}{\gamma_s^2}+O(\delta^3).
\end{equation}
Произведение \eqref{eq:C} и \eqref{eq:vsv} даёт
\begin{equation}
\label{eq:tau}
\tau=\eta_0\delta+\eta_1\delta^2,
\qquad
\eta_0=\alpha_0-\frac{1}{\gamma_s^2},
\qquad
\eta_1=\alpha_1-\frac{\alpha_0}{\gamma_s^2}+\Bigl(1+\frac{\beta_s^2}{2}\Bigr)\frac{1}{\gamma_s^2}.
\end{equation}
Через $w$ с помощью \eqref{eq:w-series}
\begin{equation}
\label{eq:tau-w}
\tau=\eta_0w+\tilde\eta_1w^2,
\qquad
\tilde\eta_1=\eta_1-\frac{\eta_0}{2\gamma_s^2}.
\end{equation}

\subsection{Система продольных уравнений}

Ускорения нет: $\varphi_s=0$ или $\pi$, и прирост энергии в зазоре $qV\sin(\varphi_s+\varphi)=qV\cos\varphi_s\sin\varphi$. За оборот
\begin{equation}
\label{eq:map}
\varphi_{k+1}-\varphi_k=2\pi h\,\tau_k,
\qquad
w_{k+1}-w_k=\frac{qV}{\beta_s^2E_s}\cos\varphi_s\sin\varphi_k .
\end{equation}
Оба равенства точные: опоздание на $T-T_s$ сдвигает фазу на $\omega_\rf(T-T_s)=2\pi h\tau$, а зазор меняет энергию, то есть $w$, на заданную величину. Деление на $T_s$ даёт систему
\begin{equation}
\label{eq:smooth}
\frac{\mathrm{d}\varphi}{\mathrm{d}t}=\omega_\rf\,\tau,
\qquad
\frac{\mathrm{d}w}{\mathrm{d}t}=K\cos\varphi_s\sin\varphi,
\qquad
K=\frac{qV}{\beta_s^2E_sT_s}.
\end{equation}
Первое уравнение описывает изменение фазы из-за запаздывания, второе --- изменение энергии в зазоре.

Устойчивой точкой $\varphi=w=0$ является
\begin{equation}
\label{eq:phis}
\cos\varphi_s=-\operatorname{sgn}\eta_0 ,
\end{equation}
малые колебания около неё имеют частоту
\begin{equation}
\label{eq:ws}
\omega_s^2=K\omega_\rf|\eta_0|,
\qquad
K\omega_\rf\cos\varphi_s=-\frac{\omega_s^2}{\eta_0}.
\end{equation}

В этой главе $\sin\varphi\approx\varphi$, и с \eqref{eq:tau-w} система принимает вид
\begin{equation}
\label{eq:system}
\frac{\mathrm{d}\varphi}{\mathrm{d}t}=\omega_\rf\bigl(\eta_0w+\tilde\eta_1w^2\bigr),
\qquad
\frac{\mathrm{d}w}{\mathrm{d}t}=K\cos\varphi_s\,\varphi.
\end{equation}

\subsection{Решение}
\label{sec:solution}

Дифференцирование второго уравнения \eqref{eq:system}, подстановка первого и \eqref{eq:ws} дают
\begin{equation}
\label{eq:ode}
\ddot w+\omega_s^2w=-\frac{\omega_s^2\tilde\eta_1}{\eta_0}\,w^2.
\end{equation}
Его решение (приложение~\ref{app:bm}):
\begin{align}
w(t)&=w_\eq+A\cos\psi+\frac{\tilde\eta_1A^2}{6\eta_0}\cos2\psi,
\qquad
w_\eq=-\frac{\tilde\eta_1A^2}{2\eta_0},
\qquad
\psi=\Omega t,
\label{eq:w-sol}\\
\varphi(t)&=\frac{\omega_\rf}{\Omega}\Bigl(\eta_0A\sin\psi+\frac{\tilde\eta_1A^2}{3}\sin2\psi\Bigr),
\label{eq:phi}\\
\frac{\Omega}{\omega_s}&=1-\frac{5}{12}\Bigl(\frac{\tilde\eta_1A}{\eta_0}\Bigr)^2.
\label{eq:Omega}
\end{align}
Амплитуда $A=(w_{\max}-w_{\min})/2$: при $\psi=0$ и $\psi=\pi$ вторая гармоника равна одной и той же величине и полуразмах не меняет. $\Omega$ --- частота синхротронных колебаний, то есть частота обхода траектории на фазовой плоскости $(\varphi,w)$.

\subsection{Усреднение и равновесный сдвиг}

$\langle f\rangle$ --- среднее по оборотам за один синхротронный период. Равновесные сдвиги энергии и импульса:
\begin{equation}
\label{eq:deq-def}
w_\eq=\langle w\rangle,
\qquad
\delta_\eq=\langle\delta\rangle .
\end{equation}
На замкнутой траектории фаза возвращается к прежнему значению, поэтому из первого уравнения \eqref{eq:smooth} точно
\begin{equation}
\label{eq:tau-avg}
\langle\tau\rangle=0 :
\end{equation}
средний период обращения частицы равен $T_s$. С \eqref{eq:tau} и \eqref{eq:tau-w}
\begin{equation}
\label{eq:deq-avg}
\delta_\eq=-\frac{\eta_1}{\eta_0}\langle\delta^2\rangle,
\qquad
w_\eq=-\frac{\tilde\eta_1}{\eta_0}\langle w^2\rangle .
\end{equation}
На решении \eqref{eq:w-sol} $\langle w^2\rangle=\langle\delta^2\rangle=A^2/2+O(A^3)$, откуда
\begin{equation}
\label{eq:deq}
\delta_\eq=-\frac{\eta_1A^2}{2\eta_0},
\qquad
w_\eq=\delta_\eq+\frac{A^2}{4\gamma_s^2}=-\frac{\tilde\eta_1A^2}{2\eta_0},
\qquad
\Delta E_\eq=\beta_s^2E_s\,w_\eq .
\end{equation}
Равенство $w_\eq=\delta_\eq+\langle\delta^2\rangle/(2\gamma_s^2)$ следует из \eqref{eq:w-series}.

Условия применимости:
\begin{equation}
\label{eq:valid}
\Bigl|\frac{\eta_1A}{\eta_0}\Bigr|\ll1,
\qquad
\hat\varphi=\frac{\omega_\rf|\eta_0|A}{\omega_s}\ll1,
\end{equation}
где $\hat\varphi$ --- полуразмах фазы.


\clearpage
\section{Магнитное кольцо: бетатронное запаздывание}
\label{sec:mag-beta}

\subsection{Запаздывание от бетатронных колебаний}

Бетатронное колебание меняет время оборота. При $w=0$ относительное запаздывание за оборот, усреднённое по быстрым бетатронным фазам,
\begin{equation}
\label{eq:Lambda}
\Lambda=\frac{\langle T-T(J=0)\rangle_\psi}{T_s}.
\end{equation}
Синхротронное движение медленное и видит только это среднее. Здесь $J=(J_x,J_y)$ --- действия, $x'=\mathrm{d}x/\mathrm{d}s$, скобки $\langle\,\rangle_\psi$ --- среднее по бетатронным фазам. $\beta_x$, $\alpha_x$, $\gamma_x=(1+\alpha_x^2)/\beta_x$ --- функции Твисса, $\rho$ --- радиус кривизны орбиты.

В магнитном поле скорость вдоль любой траектории постоянна. При $w=0$ она на траектории с данным $J$ и на траектории $J=0$ одна и та же, поэтому запаздывание по времени равно запаздыванию по пути:
\begin{equation}
\label{eq:Lambda-path}
\Lambda=\frac{\langle\Delta L_\beta\rangle_\psi}{C_s}.
\end{equation}
Равенство выполняется только потому, что скорость на обеих траекториях одинакова. $\Delta L_\beta$ --- избыток длины траектории над длиной орбиты $J=0$.

Элемент пути до второго порядка. Корень $\sqrt{(1+x/\rho)^2+x'^2+y'^2}$ после сокращения $(x/\rho)^2$ даёт
\begin{equation}
\label{eq:dl}
\mathrm{d}l=\Bigl(1+\frac{x}{\rho}+\frac{x'^2+y'^2}{2}\Bigr)\mathrm{d}s.
\end{equation}
Отклонение от орбиты $J=0$
\begin{equation}
\label{eq:xbeta}
x=x_\beta+x_{\mathrm{nl}},
\qquad
x_\beta=\sqrt{2J_x\beta_x}\cos\psi_x
\end{equation}
(для $y$ --- так же). Слагаемое $x_{\mathrm{nl}}$ --- сдвиг орбиты, квадратичный по амплитуде; его дают секступоли, кривизна и краевые поля. Среднее $\langle x_\beta\rangle_\psi=0$.

Дифференцирование \eqref{eq:xbeta} по $s$ при $\mathrm{d}\psi_x/\mathrm{d}s=1/\beta_x$ и $\alpha_x=-\beta_x'/2$ даёт
\begin{equation}
\label{eq:xp}
x_\beta'
=\sqrt{2J_x}\Bigl(\frac{\beta_x'}{2\sqrt{\beta_x}}\cos\psi_x-\frac{\sin\psi_x}{\sqrt{\beta_x}}\Bigr)
=-\sqrt{\frac{2J_x}{\beta_x}}\bigl(\sin\psi_x+\alpha_x\cos\psi_x\bigr).
\end{equation}
Усреднение $\langle\sin^2\psi_x\rangle_\psi=\langle\cos^2\psi_x\rangle_\psi=\tfrac12$, $\langle\sin\psi_x\cos\psi_x\rangle_\psi=0$ даёт $\langle x_\beta'^2\rangle_\psi=J_x\gamma_x$, и так же для $y$. В \eqref{eq:dl} среднее $\langle x'^2\rangle_\psi=\langle x_\beta'^2\rangle_\psi$: среднее $x_\beta'$ равно нулю, а $x_{\mathrm{nl}}'^2$ имеет более высокий порядок по амплитуде. Поэтому
\begin{equation}
\label{eq:Lambda-g}
\Lambda=\frac{1}{C_s}\Biggl[
\frac{J_x}{2}\oint\gamma_x\,\mathrm{d}s
+\frac{J_y}{2}\oint\gamma_y\,\mathrm{d}s
+\oint\frac{\langle x_{\mathrm{nl}}\rangle_\psi}{\rho}\,\mathrm{d}s
\Biggr].
\end{equation}

При постоянных бета-функциях $\alpha=0$, $\gamma=1/\beta$ и $\oint\mathrm{d}s/\beta=2\pi\nu$. Без $x_{\mathrm{nl}}$
\begin{equation}
\label{eq:Lambda-nu}
\Lambda=\frac{1}{C_s}\Bigl(\frac{J_x}{2}\cdot2\pi\nu_x+\frac{J_y}{2}\cdot2\pi\nu_y\Bigr)
=\frac{\pi}{C_s}\bigl(J_x\nu_x+J_y\nu_y\bigr).
\end{equation}

Та же величина через полные хроматичности $\xi_{x,y}=\mathrm{d}\nu_{x,y}/\mathrm{d}\delta$, включая секступоли; вклад $x_{\mathrm{nl}}$ входит в них автоматически:
\begin{equation}
\label{eq:Lambda-xi}
\Lambda=-\frac{2\pi}{C_s}\bigl(J_x\xi_x+J_y\xi_y\bigr).
\end{equation}
Формула следует из симплектичности отображения за оборот и верна в первом порядке по $J$~\cite{Shoji2005}. В чистом фокусирующем канале $\xi=-\nu/2$, и \eqref{eq:Lambda-xi} совпадает с \eqref{eq:Lambda-nu}.

\subsection{Система с бетатронным запаздыванием}

$\Lambda$ входит в запаздывание за оборот, то есть в первое уравнение \eqref{eq:smooth}:
\begin{equation}
\label{eq:tau-L}
\tau=\eta_0w+\tilde\eta_1w^2+\Lambda
=\eta_0\delta+\eta_1\delta^2+\Lambda.
\end{equation}
Второе выражение эквивалентно первому по \eqref{eq:tau} и \eqref{eq:tau-w}. При синхротронном движении $J_x$ и $J_y$ постоянны, поэтому $\Lambda$ постоянна. Зависимость $\Lambda$ от $w$ имеет более высокий порядок и отбрасывается.

В этой главе $\sin\varphi\approx\varphi$, и система \eqref{eq:system} заменяется на
\begin{equation}
\label{eq:system-L}
\frac{\mathrm{d}\varphi}{\mathrm{d}t}=\omega_\rf\bigl(\eta_0w+\tilde\eta_1w^2+\Lambda\bigr),
\qquad
\frac{\mathrm{d}w}{\mathrm{d}t}=K\cos\varphi_s\,\varphi.
\end{equation}

\subsection{Решение и равновесный сдвиг}

Дифференцирование второго уравнения \eqref{eq:system-L}, подстановка первого и \eqref{eq:ws} дают
\begin{equation}
\label{eq:ode-L}
\ddot w+\omega_s^2w=c_0+c_2w^2,
\qquad
c_0=-\frac{\omega_s^2\Lambda}{\eta_0},
\qquad
c_2=-\frac{\omega_s^2\tilde\eta_1}{\eta_0}.
\end{equation}
Решение (приложение~\ref{app:bm}); амплитуда $A$ имеет тот же смысл, что в \eqref{eq:w-sol}:
\begin{align}
w(t)&=w_\eq+A\cos\psi+\frac{\tilde\eta_1A^2}{6\eta_0}\cos2\psi,
\nonumber\\
w_\eq&=-\frac{1}{\eta_0}\Bigl(\frac{\tilde\eta_1A^2}{2}+\Lambda\Bigr),
\qquad
\psi=\Omega t,
\label{eq:w-sol-L}\\
\varphi(t)&=\frac{\omega_\rf}{\Omega}\Biggl[
\eta_0A\Bigl(1-\frac{2\tilde\eta_1\Lambda}{\eta_0^2}\Bigr)\sin\psi
+\frac{\tilde\eta_1A^2}{3}\sin2\psi
\Biggr],
\label{eq:phi-L}\\
\frac{\Omega}{\omega_s}&=1-\frac{\tilde\eta_1\Lambda}{\eta_0^2}-\frac{5}{12}\Bigl(\frac{\tilde\eta_1A}{\eta_0}\Bigr)^2.
\label{eq:Omega-L}
\end{align}

Из \eqref{eq:tau-avg} и \eqref{eq:tau-L}
\begin{equation}
\label{eq:deq-raw-L}
\delta_\eq=-\frac{\eta_1\langle\delta^2\rangle+\Lambda}{\eta_0}.
\end{equation}
На решении \eqref{eq:w-sol-L} $\langle\delta^2\rangle=A^2/2$, поэтому с \eqref{eq:w-series} и \eqref{eq:tau-w}
\begin{equation}
\label{eq:deq-L}
\delta_\eq=-\frac{1}{\eta_0}\Bigl(\frac{\eta_1A^2}{2}+\Lambda\Bigr),
\qquad
w_\eq=\delta_\eq+\frac{A^2}{4\gamma_s^2},
\qquad
\Delta E_\eq=\beta_s^2E_s\,w_\eq.
\end{equation}
Постоянный член \eqref{eq:w-sol-L} совпадает с этим $w_\eq$.

При $A\to0$ устойчивая точка смещается в $w=-\Lambda/\eta_0$: в удержанном порядке там $\tau=0$, а $\varphi=0$ остаётся центром колебаний фазы.

К условиям \eqref{eq:valid} добавляется
\begin{equation}
\label{eq:valid-L}
\Bigl|\frac{\tilde\eta_1\Lambda}{\eta_0^2}\Bigr|\ll1:
\end{equation}
наклон $\tau$ в смещённой устойчивой точке, $\eta_0-2\tilde\eta_1\Lambda/\eta_0$, мало отличается от $\eta_0$.


\clearpage
\section{Магнитное кольцо: второй порядок по фазе}
\label{sec:mag-phase}

\subsection{Система}

Ряд синуса обрывается на кубическом члене: $\sin\varphi\approx\varphi-\varphi^3/6$. В производной второго уравнения \eqref{eq:smooth} это замена $\cos\varphi\approx 1-\varphi^2/2$. Сначала $\Lambda=0$, и система \eqref{eq:system} принимает вид
\begin{equation}
\label{eq:system-phi}
\frac{\mathrm{d}\varphi}{\mathrm{d}t}=\omega_\rf\bigl(\eta_0 w+\tilde\eta_1 w^2\bigr),
\qquad
\frac{\mathrm{d}w}{\mathrm{d}t}=K\cos\varphi_s\Bigl(\varphi-\frac{\varphi^3}{6}\Bigr).
\end{equation}
При $\tilde\eta_1=0$ из \eqref{eq:system-phi} исключается $w$. Дифференцирование первого уравнения и подстановка второго с $\omega_\rf\eta_0 K\cos\varphi_s=-\omega_s^2$ из \eqref{eq:ws} дают
\begin{equation}
\label{eq:phi-ode}
\ddot\varphi+\omega_s^2\Bigl(\varphi-\frac{\varphi^3}{6}\Bigr)=0.
\end{equation}
Полуразмах фазы --- $\hat\varphi=\max|\varphi|$.

\subsection{Решение}

При $\tilde\eta_1=0$ решение \eqref{eq:phi-ode} получено в приложении~\ref{sec:bm-cubic}. Множитель $\operatorname{sgn}\eta_0$ уравнивает ведущий член с \eqref{eq:phi}:
\begin{align}
\varphi(t)
&=\operatorname{sgn}\eta_0\Biggl[
\hat\varphi\sin\psi
+\frac{\hat\varphi^3}{192}\bigl(\sin\psi+\sin 3\psi\bigr)
\Biggr],
\label{eq:phi-ph}\\
w(t)
&=A\Biggl[
\Bigl(1-\frac{\hat\varphi^2}{64}\Bigr)\cos\psi
+\frac{\hat\varphi^2}{64}\cos 3\psi
\Biggr],
\label{eq:w-ph}\\
\psi&=\Omega t,
\qquad
\frac{\Omega}{\omega_s}=1-\frac{\hat\varphi^2}{16}.
\label{eq:Omega-ph}
\end{align}
В ведущем порядке $\hat\varphi=\omega_\rf|\eta_0|A/\omega_s$, и \eqref{eq:phi-ph} совпадает с \eqref{eq:phi}. При $\psi=0$ и $\psi=\pi$ скобка в \eqref{eq:w-ph} равна $\pm1$, так что $A=(w_{\max}-w_{\min})/2$. Этот полуразмах связан с $\hat\varphi$ формулой
\begin{equation}
\label{eq:A-ph}
A=\frac{\omega_s\hat\varphi}{\omega_\rf|\eta_0|}\Bigl(1-\frac{\hat\varphi^2}{24}\Bigr).
\end{equation}
Для маятника $\ddot\varphi+\omega_s^2\sin\varphi=0$ интеграл энергии даёт то же соотношение точно:
\begin{equation}
\label{eq:A-exact}
A=\frac{2\omega_s}{\omega_\rf|\eta_0|}\sin\frac{\hat\varphi}{2}.
\end{equation}
Разложение \eqref{eq:A-exact} до членов $\hat\varphi^2$ совпадает с \eqref{eq:A-ph}.

При $\tilde\eta_1\neq 0$ в первом порядке по $\tilde\eta_1$ к \eqref{eq:w-ph} добавляются постоянный член и вторая гармоника \eqref{eq:w-sol}. Члены порядка $\tilde\eta_1 A^2\hat\varphi^2$ не удерживаются:
\begin{align}
w(t)
&=w_\eq
+A\Biggl[
\Bigl(1-\frac{\hat\varphi^2}{64}\Bigr)\cos\psi
+\frac{\hat\varphi^2}{64}\cos 3\psi
\Biggr]
+\frac{\tilde\eta_1 A^2}{6\eta_0}\cos 2\psi,
\label{eq:w-ph-eta}\\
\frac{\Omega}{\omega_s}
&=1-\frac{\hat\varphi^2}{16}
-\frac{5}{12}\Bigl(\frac{\tilde\eta_1 A}{\eta_0}\Bigr)^2.
\label{eq:Omega-ph-eta}
\end{align}
Постоянная $w_\eq$ определяется в \eqref{eq:deq-ph}. При $\psi=0$ и $\psi=\pi$ вторая гармоника одинакова, как в \eqref{eq:w-sol}, и $A$ остаётся полуразмахом.

\subsection{Усреднение и равновесный сдвиг}

Равенства \eqref{eq:deq-avg} точны и от второго уравнения системы не зависят: порядок по фазе входит только через $\langle w^2\rangle$. На этом порядке $\langle w^2\rangle=\langle\delta^2\rangle$. В первом порядке по $\tilde\eta_1$ среднее берётся на решении \eqref{eq:w-ph}. Гармоники $\cos\psi$ и $\cos 3\psi$ ортогональны, поэтому
\begin{equation}
\label{eq:w2-ph}
\langle w^2\rangle
=\frac{A^2}{2}\Biggl[
\Bigl(1-\frac{\hat\varphi^2}{64}\Bigr)^2
+\Bigl(\frac{\hat\varphi^2}{64}\Bigr)^2
\Biggr]
=\frac{A^2}{2}\Bigl(1-\frac{\hat\varphi^2}{32}\Bigr)+O(\hat\varphi^4).
\end{equation}
Подстановка в \eqref{eq:deq-avg} даёт
\begin{equation}
\label{eq:deq-ph}
\delta_\eq=-\frac{\eta_1 A^2}{2\eta_0}\Bigl(1-\frac{\hat\varphi^2}{32}\Bigr),
\qquad
w_\eq=-\frac{\tilde\eta_1 A^2}{2\eta_0}\Bigl(1-\frac{\hat\varphi^2}{32}\Bigr),
\qquad
\Delta E_\eq=\beta_s^2 E_s\,w_\eq .
\end{equation}

Внутри сепаратрисы для маятника $\ddot\varphi+\omega_s^2\sin\varphi=0$ при любом $\hat\varphi$
\begin{equation}
\label{eq:w2-ell}
\langle w^2\rangle
=A^2\frac{\mathrm{E}(k)/\mathrm{K}(k)-1+k^2}{k^2},
\qquad
k=\sin\frac{\hat\varphi}{2},
\end{equation}
где $\mathrm{K}$ и $\mathrm{E}$ --- полные эллиптические интегралы первого и второго рода, $k$ --- модуль. Разложение \eqref{eq:w2-ell} до членов $\hat\varphi^2$ даёт множитель $(1-\hat\varphi^2/32)$ в \eqref{eq:w2-ph}.

Поправка к $\langle w^2\rangle$ возникает потому, что $A$ есть полуразмах: кубический член рождает третью гармонику, и при фиксированном полуразмахе она перераспределяет $\langle w^2\rangle$.

\subsection{С бетатронным запаздыванием}

Первое уравнение содержит $\tau$ из \eqref{eq:tau-L}, второе --- ряд $\sin\varphi\approx\varphi-\varphi^3/6$:
\begin{equation}
\label{eq:system-phi-L}
\frac{\mathrm{d}\varphi}{\mathrm{d}t}
=\omega_\rf\bigl(\eta_0 w+\tilde\eta_1 w^2+\Lambda\bigr),
\qquad
\frac{\mathrm{d}w}{\mathrm{d}t}
=K\cos\varphi_s\Bigl(\varphi-\frac{\varphi^3}{6}\Bigr).
\end{equation}
При $\tilde\eta_1=0$ устойчивая точка лежит в $w=-\Lambda/\eta_0$. Отклонение $u=w+\Lambda/\eta_0$ удовлетворяет \eqref{eq:phi-ode}: движение около этой точки есть тот же маятник, с полуразмахом $A$ относительно смещённой точки и со связью \eqref{eq:A-ph}, \eqref{eq:A-exact}.

В первом порядке по $\tilde\eta_1$ и по $\Lambda$ величина $\langle\delta^2\rangle$ берётся из \eqref{eq:w2-ph}. Подстановка в \eqref{eq:deq-raw-L} и переход к $w$ по \eqref{eq:w-series}, \eqref{eq:tau-w} дают
\begin{equation}
\label{eq:deq-ph-L}
\begin{aligned}
\delta_\eq
&=-\frac{1}{\eta_0}\Biggl[
\frac{\eta_1 A^2}{2}\Bigl(1-\frac{\hat\varphi^2}{32}\Bigr)+\Lambda
\Biggr],
\\
w_\eq
&=\delta_\eq+\frac{A^2}{4\gamma_s^2}\Bigl(1-\frac{\hat\varphi^2}{32}\Bigr)
=-\frac{1}{\eta_0}\Biggl[
\frac{\tilde\eta_1 A^2}{2}\Bigl(1-\frac{\hat\varphi^2}{32}\Bigr)+\Lambda
\Biggr],
\\
\Delta E_\eq
&=\beta_s^2 E_s\,w_\eq .
\end{aligned}
\end{equation}
При $\Lambda=0$ это \eqref{eq:deq-ph}.

К правой части \eqref{eq:Omega-L} добавляется сдвиг $-\hat\varphi^2/16$ из \eqref{eq:Omega-ph}:
\begin{equation}
\label{eq:Omega-ph-L}
\frac{\Omega}{\omega_s}
=1-\frac{\tilde\eta_1\Lambda}{\eta_0^2}
-\frac{5}{12}\Bigl(\frac{\tilde\eta_1 A}{\eta_0}\Bigr)^2
-\frac{\hat\varphi^2}{16}.
\end{equation}

Условия применимости ряда:
\begin{equation}
\label{eq:valid-ph}
\Bigl|\frac{\eta_1 A}{\eta_0}\Bigr|\ll 1,
\qquad
\Bigl|\frac{\tilde\eta_1\Lambda}{\eta_0^2}\Bigr|\ll 1,
\qquad
\hat\varphi^2\ll 1.
\end{equation}
Формула \eqref{eq:w2-ell} верна внутри сепаратрисы при любом $\hat\varphi$.


\clearpage
\section[Электростатическое кольцо: коэффициент скольжения и равновесный сдвиг]{Электростатическое кольцо:\\ коэффициент скольжения\\ и равновесный сдвиг}
\label{sec:es-base}

В электростатическом дефлекторе часть полной энергии частицы переходит в потенциальную, поэтому скорость внутри дефлектора зависит от радиального положения. Это меняет и дисперсию, и время пролёта. В главе предложена модель потенциала дефлектора, выведены коэффициенты запаздывания электростатического кольца и решены продольные уравнения. Бетатронные колебания не учитываются: частица движется по замкнутой орбите в плоскости кольца, $y=0$.

\subsection{Постановка и потенциал дефлектора}

Поворот пучка создают только электростатические дефлекторы; вне них $1/\rho=0$ и электрический потенциал равен нулю, а фокусирующие элементы магнитные. Резонатор расположен вне дефлекторов. Во всех дефлекторах один и тот же показатель поля $n$; в плоскости кольца
\begin{equation}
\label{eq:Ex}
E_x=E_0\Bigl(1-n\frac{x}{\rho}+O(x^2)\Bigr),
\end{equation}
где $E_0$ --- поле на референсной орбите, $x$ --- радиальное отклонение, $y$ --- вертикальное. В дефлекторе продольный масштабный множитель $h_s=1+x/\rho$.

Центростремительная сила на референсной орбите задаёт
\begin{equation}
\label{eq:E0}
qE_0=-\frac{p_sv_s}{\rho}.
\end{equation}
Потенциал $\Phi$ отсчитывается от референсной орбиты и измеряется в единицах $p_sv_s=\beta_s^2E_s$:
\begin{equation}
\label{eq:psi-def}
\chi=\frac{q\Phi}{p_sv_s}.
\end{equation}
Интегрирование $E_x=-\partial\Phi/\partial x$ с \eqref{eq:Ex}, \eqref{eq:E0} даёт члены по $x$, а уравнение Лапласа в координатах дефлектора (поле не зависит от $s$)
\begin{equation}
\label{eq:laplace}
\frac{\partial^2\chi}{\partial x^2}
+\frac{\partial^2\chi}{\partial y^2}
+\frac{1}{\rho+x}\frac{\partial\chi}{\partial x}=0
\end{equation}
фиксирует коэффициент при $y^2$:
\begin{equation}
\label{eq:psi}
\chi
=\frac{x}{\rho}
-\frac{n}{2}\frac{x^2}{\rho^2}
+\frac{n-1}{2}\frac{y^2}{\rho^2}
+O\bigl((x/\rho)^3,(y/\rho)^3\bigr).
\end{equation}
На внешнем радиусе $x>0$ потенциальная энергия частицы положительна. Член с $y^2$ на орбите в плоскости кольца не участвует; он нужен для вертикальных бетатронных колебаний.

\subsection{Полная и кинетическая энергия, скорость в дефлекторе}

Поле дефлектора статическое, поэтому между прохождениями резонатора сохраняется \emph{полная} энергия $\gamma mc^2+q\Phi$. \emph{Кинетическая} энергия $\gamma mc^2$ внутри дефлектора меняется вместе с потенциалом: при входе на внешний радиус часть её переходит в потенциальную, при выходе возвращается. Лоренц-фактор $\gamma$ всюду ниже --- кинетический.

Переменная $\delta$ определяется по \eqref{eq:delta-def} через импульс вне дефлекторов, где $\Phi=0$. Тем самым $\delta$ --- метка полной энергии, одинаковая вдоль всего оборота. Обозначим через $\gamma_{\mathrm{out}}(\delta)$ лоренц-фактор вне дефлекторов; он даётся рядом \eqref{eq:gamma-delta}. Сохранение полной энергии, $\gamma mc^2+q\Phi=\gamma_{\mathrm{out}}mc^2$, даёт локальный лоренц-фактор
\begin{equation}
\label{eq:gamma-loc}
\gamma=\gamma_{\mathrm{out}}(\delta)-\beta_s^2\gamma_s\chi
=\gamma_s+\beta_s^2\gamma_s(\delta-\chi)
+\frac{\beta_s^2}{2\gamma_s}\delta^2
+O(\delta^3).
\end{equation}
Член с $\chi$ точен, так как потенциал входит в полную энергию линейно.

Скорость в данной точке удобно записать через локальный импульс. Определим $\delta_{\mathrm{loc}}$ как значение $\delta$, при котором ряд \eqref{eq:gamma-delta} даёт локальный $\gamma$. Из \eqref{eq:gamma-loc}
\begin{equation}
\label{eq:delta-loc}
\delta_{\mathrm{loc}}
=\delta-\chi+\frac{2\delta\chi-\chi^2}{2\gamma_s^2}+O(3),
\end{equation}
где $O(3)$ --- члены суммарной третьей степени по $\delta$ и $\chi$. Это не новая динамическая переменная, а значение локального импульса в данной точке. Ряд \eqref{eq:vsv} в точке $\delta_{\mathrm{loc}}$ даёт
\begin{equation}
\label{eq:vsv-es}
\frac{v_s}{v}
=1-\frac{\delta-\chi}{\gamma_s^2}
+\frac{3(\delta-\chi)^2}{2\gamma_s^2}
-\frac{\delta^2}{2\gamma_s^4}
+O(3).
\end{equation}
На внешнем радиусе частица медленнее, чем в свободном промежутке при том же $\delta$.

Резонатор стоит при $\chi=0$, поэтому отображение \eqref{eq:map} и система \eqref{eq:smooth} сохраняют свой вид. От дефлекторов зависит только функция $\tau(\delta)$.

\subsection{Дисперсия в дефлекторе}
\label{sec:es-D}

Кривизну траектории задаёт произведение $pv$, то есть $1/R=-qE_x/(pv)$. Из \eqref{eq:gamma-loc} при $\chi=x/\rho$ и $pv=mc^2(\gamma-1/\gamma)$ в линейном порядке
\begin{equation}
\label{eq:cmp-pv}
\frac{pv}{p_sv_s}
=1+\Bigl(1+\frac{1}{\gamma_s^2}\Bigr)
\Bigl(\delta-\frac{x}{\rho}\Bigr).
\end{equation}
Подстановка в уравнение орбиты в сопровождающей системе, $x''+x/\rho^2=1/\rho-1/R$, даёт для замкнутой орбиты $x=D_1\delta$ внутри дефлектора
\begin{equation}
\label{eq:cmp-D-es}
D_1''+\frac{2-n+1/\gamma_s^2}{\rho^2}\,D_1
=\frac{1+1/\gamma_s^2}{\rho}.
\end{equation}
В жёсткости единица --- центробежный член, $-n$ --- градиент поля, а $1+1/\gamma_s^2$ --- уменьшение $pv$ на внешнем радиусе из-за потенциала: там то же поле поворачивает траекторию сильнее. Источник $(1+1/\gamma_s^2)/\rho$ вместо магнитного $1/\rho$ отражает то, что кривизну задаёт $pv$, а не $p$~\cite{Mane2008,Mane2012}. Вне дефлекторов уравнение переходит в $D_1''+K_1D_1=0$. На границах дефлектора $D_1$ и $D_1'$ непрерывны: линейной краевой линзы при перпендикулярном торце нет. Во втором порядке непрерывны канонические импульсы, а связь $p_x$ с наклоном внутри и снаружи разная, $p_x=x'-(2/\rho)xx'$ внутри и $p_x=x'$ снаружи. Поэтому $D_2$ непрерывно, а $D_2'$ по разные стороны границы различается: на входе $D_2^{\prime\mathrm{in}}-D_2^{\prime\mathrm{out}}=2D_1D_1'/\rho$, на выходе --- с обратным знаком. Это граничное условие для $D_2$; среднее $\langle D_2/\rho\rangle_C$ входит в $\eta_1^{\mathrm{E}}$ \eqref{eq:eta1-es}. Протяжённое краевое поле электродов здесь не учитывается.

\subsection{Запаздывание за оборот: коэффициенты $\eta_0^{\mathrm{E}}$ и $\eta_1^{\mathrm{E}}$}

Замкнутая орбита
\begin{equation}
\label{eq:x-co}
x=D_1\delta+D_2\delta^2
\end{equation}
определена по внешнему импульсу $\delta$. Период обращения $T=\oint\mathrm{d}l/v$, поэтому
\begin{equation}
\label{eq:T-oint}
1+\tau=\frac{v_s}{C_s}\oint\frac{\mathrm{d}l}{v}
=\Bigl\langle\frac{\mathrm{d}l}{\mathrm{d}s}\,\frac{v_s}{v}\Bigr\rangle_C,
\end{equation}
где
\begin{equation}
\label{eq:avg-C}
\langle f\rangle_C=\frac{1}{C_s}\oint f\,\mathrm{d}s
\end{equation}
--- среднее вдоль референсной орбиты (в отличие от среднего $\langle f\rangle$ по синхротронному периоду). Элемент пути при $y=0$, как и в \eqref{eq:dl},
\begin{equation}
\label{eq:dl-es}
\frac{\mathrm{d}l}{\mathrm{d}s}
=1+\frac{x}{\rho}+\frac{x'^2}{2}+O(x^3).
\end{equation}
На орбите \eqref{eq:x-co} потенциал
\begin{equation}
\label{eq:psi-co}
\chi
=\frac{D_1}{\rho}\,\delta
+\Bigl[\frac{D_2}{\rho}-\frac{n}{2}\Bigl(\frac{D_1}{\rho}\Bigr)^2\Bigr]\delta^2
+O(\delta^3),
\end{equation}
и произведение \eqref{eq:dl-es} на \eqref{eq:vsv-es} равно
\begin{equation}
\label{eq:integrand}
\begin{aligned}
\frac{\mathrm{d}l}{\mathrm{d}s}\frac{v_s}{v}
&=1+\Bigl[
\Bigl(1+\frac{1}{\gamma_s^2}\Bigr)\frac{D_1}{\rho}
-\frac{1}{\gamma_s^2}
\Bigr]\delta
\\
&\quad+\Bigl[
\frac{D_1'^2}{2}
+\Bigl(1+\frac{1}{\gamma_s^2}\Bigr)\frac{D_2}{\rho}
+\frac{5-n}{2\gamma_s^2}\Bigl(\frac{D_1}{\rho}\Bigr)^2
\\
&\qquad\qquad
-\frac{4}{\gamma_s^2}\frac{D_1}{\rho}
+\frac{3\gamma_s^2-1}{2\gamma_s^4}
\Bigr]\delta^2
+O(\delta^3).
\end{aligned}
\end{equation}
Коэффициент $(5-n)/2$ собирается из квадратичного члена потенциала ($-n/2$), из произведения удлинения дуги на линейное замедление ($+1$) и из квадратичного члена ряда скорости ($+3/2$). С геометрическим коэффициентом уплотнения
\begin{equation}
\label{eq:alpha0-es}
\alpha_0=\Bigl\langle\frac{D_1}{\rho}\Bigr\rangle_C,
\end{equation}
вычисленным на дисперсии \eqref{eq:cmp-D-es}, запаздывание имеет вид
\begin{equation}
\label{eq:tau-es-d}
\tau=\eta_0^{\mathrm{E}}\delta+\eta_1^{\mathrm{E}}\delta^2,
\end{equation}
\begin{align}
\eta_0^{\mathrm{E}}
&=\Bigl(1+\frac{1}{\gamma_s^2}\Bigr)\alpha_0-\frac{1}{\gamma_s^2},
\label{eq:eta0-es}\\
\eta_1^{\mathrm{E}}
&=\frac12\langle D_1'^2\rangle_C
+\Bigl(1+\frac{1}{\gamma_s^2}\Bigr)
\Bigl\langle\frac{D_2}{\rho}\Bigr\rangle_C
+\frac{5-n}{2\gamma_s^2}
\Bigl\langle\Bigl(\frac{D_1}{\rho}\Bigr)^2\Bigr\rangle_C
-\frac{4\alpha_0}{\gamma_s^2}
+\frac{3\gamma_s^2-1}{2\gamma_s^4}.
\label{eq:eta1-es}
\end{align}
Индекс $\mathrm{E}$ отмечает коэффициенты электростатического кольца; ниже $\eta_0^{\mathrm{E}}\neq 0$.

\paragraph{Два источника отличий от магнитного кольца.}
Поправки порядка $1/\gamma_s^2$ входят в $\eta^{\mathrm{E}}$ из двух независимых мест. Первое --- орбита: дисперсия удовлетворяет \eqref{eq:cmp-D-es}, а не \eqref{eq:D1}, поэтому $\alpha_0$ отличается от магнитного кольца с теми же квадруполями. Второе --- время пролёта по уже построенной орбите: на смещённом радиусе частица медленнее. В
\begin{equation}
\label{eq:eta0-es-split}
\eta_0^{\mathrm{E}}
=\alpha_0+\frac{\alpha_0}{\gamma_s^2}-\frac{1}{\gamma_s^2}
\end{equation}
слагаемое $\alpha_0$ --- удлинение орбиты (первое место), $\alpha_0/\gamma_s^2$ --- среднее замедление на потенциале (второе место), $-1/\gamma_s^2$ --- зависимость скорости от $\delta$ при нулевом потенциале, общая с магнитным кольцом. В $\eta_1^{\mathrm{E}}$ средние $\frac12\langle D_1'^2\rangle_C$ и $\langle D_2/\rho\rangle_C$ относятся к орбите, а множитель $1/\gamma_s^2$ при $\langle D_2/\rho\rangle_C$, член с $(5-n)$ и член $-4\alpha_0/\gamma_s^2$ --- к времени пролёта.

В нерелятивистском пределе $\eta_0^{\mathrm{E}}\to 2\alpha_0-1$, что согласуется с измерениями на электростатическом кольце CSR~\cite{Grieser2017} (там коэффициент определён через частоту обращения и имеет противоположный знак).

\subsection{Система уравнений и равновесный сдвиг}

При $\sin\varphi\approx\varphi$
\begin{equation}
\label{eq:system-es}
\left\{
\begin{aligned}
\dot\varphi
&=\omega_\rf\bigl(\eta_0^{\mathrm{E}}\delta+\eta_1^{\mathrm{E}}\delta^2\bigr),\\
\dot\delta
&=K\cos\varphi_s\,\varphi\Bigl(1-\frac{\delta}{\gamma_s^2}\Bigr).
\end{aligned}
\right.
\end{equation}
Это система \eqref{eq:system} с заменой $\eta_0\to\eta_0^{\mathrm{E}}$, $\eta_1\to\eta_1^{\mathrm{E}}$. Устойчивая фаза и частота малых колебаний даются \eqref{eq:phis}, \eqref{eq:ws}, а решение --- \eqref{eq:delta-sol} с той же заменой. Условие $\langle\tau\rangle=0$ даёт
\begin{equation}
\label{eq:deq-es}
\delta_\eq=-\frac{\eta_1^{\mathrm{E}}\delta_m^2}{2\eta_0^{\mathrm{E}}}.
\end{equation}

Вне дефлекторов лоренц-фактор связан с $\delta$ рядом \eqref{eq:gamma-delta}, поэтому
\begin{equation}
\label{eq:gamma-eff-es}
\gamma_{\mathrm{eff}}
=\gamma_s+\beta_s^2\gamma_s\,\delta_\eq
+\frac{\beta_s^2\delta_m^2}{4\gamma_s}
+O(\delta_m^4),
\end{equation}
\begin{equation}
\label{eq:dE-es}
\Delta E_\eq
=\beta_s^2E_s\Bigl(
\delta_\eq+\frac{\delta_m^2}{4\gamma_s^2}
\Bigr)
+O(\delta_m^4 E_s).
\end{equation}
Формула \eqref{eq:dE-es} даёт сдвиг средней полной энергии; вне дефлекторов он совпадает со сдвигом кинетической. Лоренц-фактор внутри дефлектора рассмотрен в разделе~\ref{sec:es-gamma-in}.

Условия применимости --- \eqref{eq:valid} с коэффициентами этого кольца и дополнительно
\begin{equation}
\label{eq:valid-es}
\Bigl|\frac{D_1}{\rho}\Bigr|\delta_m\ll 1,
\end{equation}
т.\,е. малость потенциала в дефлекторах. Это условие существенно: например, при $n=2$ и $D_1''=0$ из \eqref{eq:cmp-D-es} следует $D_1/\rho=\gamma_s^2+1$.

\subsection{Кубический член ВЧ-напряжения}

Запаздывание \eqref{eq:tau-es-d} имеет тот же вид, что в магнитном кольце, а резонатор стоит при нулевом потенциале. Поэтому результаты главы~\ref{sec:mag-phase}, в частности \eqref{eq:deq-ph}, переносятся с заменой $\eta_0\to\eta_0^{\mathrm{E}}$, $\eta_1\to\eta_1^{\mathrm{E}}$.


\clearpage
\section[Электростатическое кольцо: бетатронное запаздывание]{Электростатическое кольцо:\\ бетатронное запаздывание}
\label{sec:es-beta}

\subsection{Запаздывание}

$\Lambda$ определяется как в \eqref{eq:Lambda}. В дефлекторе при $w=0$ выполнено $u=-\psi$, и скорость меняется вдоль бетатронной траектории, поэтому запаздывание по времени не сводится к запаздыванию по пути. Подынтегральное выражение \eqref{eq:tau-es} с элементом пути \eqref{eq:dl} и с $v_s/v$ из \eqref{eq:vsv-es} есть
\[
\Bigl(1+\frac{x}{\rho}+\frac{x'^2+y'^2}{2}\Bigr)
\Bigl(1+\frac{\psi}{\gamma_s^2}+\frac{3\psi^2}{2\gamma_s^2}\Bigr),
\]
$x=x_\beta+x_{\mathrm{nl}}$, $y=y_\beta$. Квадратичные по $x_\beta$ члены из $(x/\rho)(\psi/\gamma_s^2)$, из $-(n/2)(x/\rho)^2/\gamma_s^2$ и из $3(x/\rho)^2/(2\gamma_s^2)$ в сумме равны $(5-n)x_\beta^2/(2\gamma_s^2\rho^2)$. По $y$ остаётся $(n-1)y^2/(2\gamma_s^2\rho^2)$, линейный по $x_{\mathrm{nl}}$ член равен $(1+1/\gamma_s^2)x_{\mathrm{nl}}/\rho$. Слагаемое $(x'^2+y'^2)/2$ сохраняется. Линейные по $x_\beta$ члены в среднем $\langle\,\rangle_\psi$ равны нулю, поэтому
\begin{equation}
\label{eq:Lambda-es}
\begin{aligned}
\Lambda=\frac{1}{C_s}\oint\Biggl[&
\Bigl(1+\frac{1}{\gamma_s^2}\Bigr)\frac{\langle x_{\mathrm{nl}}\rangle_\psi}{\rho}
+\frac{\langle x'^2+y'^2\rangle_\psi}{2}
\\
&+\frac{5-n}{2\gamma_s^2}\,\frac{\langle x_\beta^2\rangle_\psi}{\rho^2}
+\frac{n-1}{2\gamma_s^2}\,\frac{\langle y_\beta^2\rangle_\psi}{\rho^2}
\Biggr]\mathrm{d}s.
\end{aligned}
\end{equation}
Среднее $x_\beta'$ равно нулю, а $x_{\mathrm{nl}}'^2$ имеет более высокий порядок, так что $\langle x'^2+y'^2\rangle_\psi=\langle x_\beta'^2\rangle_\psi+\langle y_\beta'^2\rangle_\psi$. Подстановка $\langle x_\beta'^2\rangle_\psi=J_x\gamma_x$, $\langle y_\beta'^2\rangle_\psi=J_y\gamma_y$, $\langle x_\beta^2\rangle_\psi=J_x\beta_x$, $\langle y_\beta^2\rangle_\psi=J_y\beta_y$ даёт
\begin{equation}
\label{eq:Lambda-es-J}
\begin{aligned}
\Lambda=\frac{1}{C_s}\Biggl[&
\frac{J_x}{2}\oint\Biggl(\gamma_x+\frac{(5-n)\beta_x}{\gamma_s^2\rho^2}\Biggr)\mathrm{d}s
\\
&+\frac{J_y}{2}\oint\Biggl(\gamma_y+\frac{(n-1)\beta_y}{\gamma_s^2\rho^2}\Biggr)\mathrm{d}s
\\
&+\Biggl(1+\frac{1}{\gamma_s^2}\Biggr)\oint\frac{\langle x_{\mathrm{nl}}\rangle_\psi}{\rho}\,\mathrm{d}s
\Biggr].
\end{aligned}
\end{equation}
Вне дефлекторов $1/\rho=0$. Соотношение \eqref{eq:Lambda-xi} выполняется и здесь, с хроматичностями $\xi_{x,y}=\mathrm{d}\nu_{x,y}/\mathrm{d}\delta$ электростатического кольца.

\subsection{Система и равновесный сдвиг}

$\tau=\eta_0w+\tilde\eta_1w^2+\Lambda$, где $\eta_0$ и $\eta_1$ берутся из \eqref{eq:eta0-es}, \eqref{eq:eta1-es}, $\tilde\eta_1=\eta_1-\eta_0/(2\gamma_s^2)$, а $\Lambda$ --- из \eqref{eq:Lambda-es}. Система \eqref{eq:system-L}, её решение и усреднение главы~\ref{sec:mag-beta} сохраняются:
\begin{equation}
\label{eq:deq-es-L}
\begin{aligned}
\delta_\eq&=-\frac{1}{\eta_0}\Bigl(\frac{\eta_1A^2}{2}+\Lambda\Bigr),
\\
w_\eq&=-\frac{1}{\eta_0}\Bigl(\frac{\tilde\eta_1A^2}{2}+\Lambda\Bigr),
\\
\Delta E_\eq&=\beta_s^2E_s\,w_\eq.
\end{aligned}
\end{equation}
Условия применимости --- \eqref{eq:valid-L}. Во втором порядке по фазе выполняется \eqref{eq:deq-ph-L} с этими $\eta$ и $\Lambda$.

\subsection{Кинетическая энергия внутри дефлектора}

Из \eqref{eq:gamma-loc} точно
\begin{equation}
\label{eq:kin-loc}
\frac{\gamma-\gamma_s}{\beta_s^2\gamma_s}=w-\psi.
\end{equation}
Усреднение по синхротронным и бетатронным колебаниям в точке $s$ дефлектора даёт $\langle\gamma-\gamma_s\rangle/(\beta_s^2\gamma_s)=w_\eq-\langle\psi\rangle(s)$. При $x=D\delta+D_2\delta^2+x_\beta+x_{\mathrm{nl}}$ и $y=y_\beta$, с $\langle\delta\rangle=\delta_\eq$ и $\langle\delta^2\rangle=A^2/2$, члены $\delta\cdot x_\beta$ в среднем равны нулю, поскольку фазы независимы. Из \eqref{eq:psi}
\begin{equation}
\label{eq:psi-avg}
\begin{aligned}
\langle\psi\rangle
&=\frac{D}{\rho}\,\delta_\eq
+\Biggl[\frac{D_2}{\rho}-\frac{n}{2}\Biggl(\frac{D}{\rho}\Biggr)^2\Biggr]\frac{A^2}{2}
+\frac{\langle x_{\mathrm{nl}}\rangle_\psi}{\rho}
\\
&\quad-\frac{n}{2}\,\frac{J_x\beta_x}{\rho^2}
+\frac{n-1}{2}\,\frac{J_y\beta_y}{\rho^2}.
\end{aligned}
\end{equation}
Средний сдвиг кинетической энергии в точке $s$
\begin{equation}
\label{eq:dK}
\Delta K_\eq(s)=\beta_s^2E_s\bigl[w_\eq-\langle\psi\rangle(s)\bigr].
\end{equation}
Вне дефлекторов $\psi=0$ и $\Delta K_\eq=\Delta E_\eq$.


\clearpage
\section{Сравнение}
\label{sec:compare}

Индекс $\mathrm{M}$ отмечает магнитное кольцо глав~\ref{sec:mag-base}--\ref{sec:mag-phase}, индекс $\mathrm{E}$ --- электростатическое кольцо глав~\ref{sec:es-base} и~\ref{sec:es-beta}. В обоих кольцах равновесный сдвиг имеет один вид. Без бетатронного запаздывания и в первом порядке по фазе это \eqref{eq:deq} и \eqref{eq:deq-es}. С бетатронным запаздыванием --- \eqref{eq:deq-L} и \eqref{eq:deq-es-L}. Во втором порядке по фазе без $\Lambda$ результат \eqref{eq:deq-ph} переносится на электростатическое кольцо с заменой $\eta_0$ и $\eta_1$ на $\eta^{\mathrm{E}}$. Этот перенос записан в подразделе второго порядка по фазе главы~\ref{sec:es-base}. С запаздыванием $\Lambda$ магнитный результат --- три строки \eqref{eq:deq-ph-L}. Для электростатического кольца в них стоят $\eta^{\mathrm{E}}$ и $\Lambda^{\mathrm{E}}$. Множитель $(1-\hat\varphi^2/32)$ умножает $\langle\delta^2\rangle$ и не умножает $\Lambda$.

Кольца различаются функциями $D_1$, $D_2$ и оптикой, по которым считаются $\eta_0$, $\eta_1$ и $\Lambda$. Внутри дефлектора скорость зависит от потенциала. Ниже $\eta_0^{\mathrm{E}}$ и $\eta_1^{\mathrm{E}}$ записаны как выражения \eqref{eq:tau} на орбите электростатического кольца плюс поправка. Поправка собрана из членов, пропорциональных $1/\gamma_s^2$. Их источник --- скорость в потенциале, ряд \eqref{eq:vsv-es}. Сами $D_1$ и функции Твисса у двух колец разные, потому что поперечные уравнения разные.

\subsection{Дисперсия и длина орбиты}

Функции $D_1(s)$ и $D_2(s)$ определены замкнутой орбитой \eqref{eq:xbeta}, \eqref{eq:x-co}: при нулевой бетатронной амплитуде $x=D_1\delta+D_2\delta^2$. Оба коэффициента периодичны с периодом кольца. В электростатическом кольце аргумент $\delta$ --- относительное отклонение модуля импульса вне дефлекторов.

На этой орбите элемент пути \eqref{eq:dl} после интегрирования по кольцу даёт длину \eqref{eq:C}. Подстановка $x=D_1\delta+D_2\delta^2$ и $x'=D_1'\delta+O(\delta^2)$ оставляет при $\delta$ среднее $D_1/\rho$, а при $\delta^2$ --- среднее $D_2/\rho$ и половину среднего $D_1'^2$:
\begin{equation}
\label{eq:cmp-alpha}
\alpha_0=\Biggl\langle\frac{D_1}{\rho}\Biggr\rangle_C,
\qquad
\alpha_1=\Biggl\langle\frac{D_2}{\rho}\Biggr\rangle_C+\frac12\langle D_1'^2\rangle_C.
\end{equation}
Среднее $\langle\,\rangle_C$ определено в \eqref{eq:avg-C}. Множитель $1/\rho$ стоит потому, что дуга радиуса $\rho+x$ длиннее опорного элемента $\mathrm{d}s$ в $(\rho+x)/\rho$ раз: относительное удлинение в данной точке есть $x/\rho$. На прямолинейном промежутке $1/\rho=0$, и радиальное смещение длину не меняет. Наклон $D_1'$ удлиняет путь и там, где поворота нет. Та же связь получается из \eqref{eq:dl-es}: квадраты $(x/\rho)^2$ в элементе пути уже сократились в главах~\ref{sec:mag-beta} и~\ref{sec:es-base}.

Формула \eqref{eq:cmp-alpha} геометрическая. Какая орбита отвечает данному $\delta$, определяют уравнения движения. Дальше $\alpha_0^{\mathrm{M}}$, $\alpha_1^{\mathrm{M}}$ --- средние \eqref{eq:cmp-alpha} на дисперсии магнитного кольца, $\alpha_0^{\mathrm{E}}$, $\alpha_1^{\mathrm{E}}$ --- на дисперсии электростатического.

В магнитном кольце $D_1$ --- периодическое решение \eqref{eq:D1}:
\begin{equation}
\label{eq:cmp-D-mag}
D_1''+K_x D_1=\frac{1}{\rho}.
\end{equation}
Жёсткость $K_x$ --- коэффициент при $x$ в \eqref{eq:hill}. В секторном диполе $K_x=K_1+1/\rho^2$. На прямолинейном промежутке правая часть равна нулю.

В дефлекторе $D_1$ удовлетворяет уравнению \eqref{eq:cmp-D-es} из раздела~\ref{sec:es-D}. Правая часть там равна $(1+1/\gamma_s^2)/\rho$, коэффициент при $D_1$ равен $(2-n+1/\gamma_s^2)/\rho^2$. Слагаемые $1/\gamma_s^2$ происходят из зависимости $pv$ от потенциала. На внешнем радиусе $\psi>0$, кинетическая энергия ниже, и то же поле поворачивает траекторию сильнее. В \eqref{eq:cmp-D-mag} на этих местах стоят $1/\rho$ и $K_x$. Вне дефлекторов $1/\rho=0$ и $\psi=0$, и \eqref{eq:cmp-D-es} переходит в $D_1''+K_1 D_1=0$.

Коэффициент $D_2$ удовлетворяет тому же линейному оператору, что и $D_1$. Правая часть квадратична по $D_1$ и $D_1'$. Её дают квадрат смещения в геометрии и в поле, а также секступоли. Главы~\ref{sec:mag-base}--\ref{sec:es-beta} это уравнение не выписывают. В $\eta_0$ и $\eta_1$ входят не сами функции, а три средних \eqref{eq:cmp-alpha}.

\subsection{Интеграл бетатронного запаздывания}

Магнитное запаздывание $\Lambda^{\mathrm{M}}$ через наклоны и сдвиг центра есть \eqref{eq:Lambda-slopes}. При тонком повороте та же величина записана в \eqref{eq:Lambda-g}. Она считается по функциям Твисса \eqref{eq:twiss-x} и по $x_{\mathrm{nl}}$. Средние наклонов даны в \eqref{eq:xp2}: $\langle x_\beta'^2\rangle_\psi=J_x\gamma_x$ и $\langle y_\beta'^2\rangle_\psi=J_y\gamma_y$. Здесь $\gamma_x$ и $\gamma_y$ --- функции Твисса, не лоренц-фактор.

Электростатический интеграл \eqref{eq:Lambda-es} использует те же наклоны и тот же $x_{\mathrm{nl}}$. К ним добавляются члены, пропорциональные $1/\gamma_s^2$. Коэффициент при $\langle x_\beta^2\rangle_\psi/\rho^2$ равен $(5-n)/(2\gamma_s^2)$, при $\langle y_\beta^2\rangle_\psi/\rho^2$ --- $(n-1)/(2\gamma_s^2)$, при $x_{\mathrm{nl}}/\rho$ --- $1+1/\gamma_s^2$. Вне дефлекторов $1/\rho=0$, и в интеграле остаётся наклон. Слагаемое $x_{\mathrm{nl}}$ уже есть среднее по бетатронной фазе, как в главе~\ref{sec:mag-beta}. Подстановка $\langle x_\beta^2\rangle_\psi=J_x\beta_x$, $\langle y_\beta^2\rangle_\psi=J_y\beta_y$ и \eqref{eq:xp2} переводит \eqref{eq:Lambda-es} в \eqref{eq:Lambda-es-J}.

Сдвиг $x_{\mathrm{nl}}$ квадратичен по амплитуде. Колебание на бетатронной частоте собрано в $x_\beta$. В магнитном кольце этот сдвиг дают секступоли, кривизна и краевые поля. В электростатическом кольце к ним добавляется зависимость поля от радиуса. При $E_r\propto r^{-n}$ глава~\ref{sec:es-beta} получает $x_{\mathrm{nl}}\neq 0$ при нулевых секступолях. Общей замкнутой формулы для произвольной решётки нет.

Если $\beta_x$ и $\beta_y$ постоянны, то $\alpha_x=\alpha_y=0$ и $\gamma_x=1/\beta_x$. По \eqref{eq:nu-def} выполнено $\oint\gamma_x\,\mathrm{d}s=2\pi\nu_x$. Вклад наклона в \eqref{eq:Lambda-es} равен $\pi J_x\nu_x/C_s$, и так же в вертикальной плоскости. Сумма этих двух вкладов есть \eqref{eq:Lambda-nu}.

Члены с $1/\rho^2$ берутся из \eqref{eq:Lambda-es} с коэффициентами $(5-n)/(2\gamma_s^2)$ и $(n-1)/(2\gamma_s^2)$. При одном и том же $\rho$ во всех дефлекторах и при полном угле поворота $2\pi$ длина дефлекторов равна $2\pi\rho$, поэтому $\oint\mathrm{d}s/\rho^2=2\pi/\rho$. В том же приближении $\beta_x=C_s/(2\pi\nu_x)$ и $\beta_y=C_s/(2\pi\nu_y)$. Подстановка $\langle x_\beta^2\rangle_\psi=J_x\beta_x$ и $\langle y_\beta^2\rangle_\psi=J_y\beta_y$ даёт оценку
\begin{equation}
\label{eq:cmp-smooth}
\begin{aligned}
\Lambda^{\mathrm{E}}
&\approx\frac{\pi}{C_s}\bigl(J_x\nu_x+J_y\nu_y\bigr)
+\frac{J_x(5-n)}{2\gamma_s^2\nu_x\rho}
+\frac{J_y(n-1)}{2\gamma_s^2\nu_y\rho}\\
&\quad+\Biggl(1+\frac{1}{\gamma_s^2}\Biggr)
\frac{1}{C_s}\oint\frac{x_{\mathrm{nl}}}{\rho}\,\mathrm{d}s.
\end{aligned}
\end{equation}
Первое слагаемое --- вклад наклонов, \eqref{eq:Lambda-nu}. Два следующих --- потенциал при постоянных бета-функциях и при общем радиусе дефлекторов. Слагаемое с $x_{\mathrm{nl}}$ перенесено из \eqref{eq:Lambda-es} вместе с множителем $1+1/\gamma_s^2$: постоянная бета-функция этот сдвиг не задаёт. На общей решётке \eqref{eq:cmp-smooth} --- оценка. Рабочий интеграл --- \eqref{eq:Lambda-es}. Через хроматичность то же запаздывание записано в \eqref{eq:Lambda-es-xi}.

\subsection{Хроматичность}

В обоих кольцах $\Lambda=-2\pi(J_x\xi_x+J_y\xi_y)/C_s$ с полной хроматичностью своего кольца. Для магнитного кольца это \eqref{eq:Lambda-xi}. Для электростатического --- \eqref{eq:Lambda-es-xi}, с $\xi_x^{\mathrm{E}}$ и $\xi_y^{\mathrm{E}}$. Производные берутся при $\delta=0$ и при нулевых действиях. В электростатическом кольце аргумент $\delta$ --- модуль импульса вне дефлекторов, \eqref{eq:delta-def}.

В магнитном кольце \eqref{eq:Lambda-xi} выведено через удлинение пути для тонких поворотов. Сравниваются \eqref{eq:Lambda-g}, \eqref{eq:xi-x} и \eqref{eq:xi-y} при $K_x=K_1$. Для секторных диполей эти явные интегралы нужно дополнить. Само равенство с полными $\xi_x$ и $\xi_y$ проверено численно на ячейке FODO с секторными диполями и тонкими квадруполями. Относительная точность --- $10^{-4}$ в плоскости орбиты и $10^{-8}$ в вертикальной плоскости. Краевое поле в этот счёт не входило.

В электростатическом кольце почленное сравнение \eqref{eq:Lambda-es-J} с \eqref{eq:xi-es} до \eqref{eq:Lambda-es-xi} не доведено. Вертикальная явная часть \eqref{eq:dKy-expl} согласована с наклоном и с членом $(n-1)$. Горизонтальный остаток \eqref{eq:dKx-gap}, члены с наклонами в уравнении для $x$, расхождение $2(n-1)/\rho^3$ и стык дефлектора с прямолинейным промежутком остаются вне этой выкладки. Равенство \eqref{eq:Lambda-es-xi} проверено численно на периодической ячейке главы~\ref{sec:es-beta}. В ячейке секторные дефлекторы, прямолинейные промежутки и тонкие магнитные квадруполи. Отдельный счёт выполнен с тонкими секступолями.

Кольца различаются составом $\xi$. В магнитном кольце градиент делится на $1+\delta$, а секступоль на орбите $x=D_1\delta$ добавляет $K_2 D_1$, как в \eqref{eq:xi-x} и \eqref{eq:xi-y}. В электростатическом кольце на магнитных элементах остаётся \eqref{eq:dK-mag}. Тело дефлектора даёт жёсткости \eqref{eq:Kx-es} и \eqref{eq:Ky-es}. Показатель $n$ переносит фокусировку из плоскости орбиты в вертикальную плоскость: в $K_x$ входит $2-n$, в $K_y$ входит $n-1$. Множитель $1/\gamma_s^2$ в $K_x$ --- замедление на внешнем радиусе, где $pv$ меньше. В магнитном диполе этого вклада нет.

Если $\xi_x=\xi_y=0$, правая часть \eqref{eq:Lambda-xi} равна нулю, и $\Lambda^{\mathrm{M}}=0$. Это относится к тонким секступолям и к первому порядку по действию. В электростатическом кольце обнуление $\xi_x^{\mathrm{E}}$ и $\xi_y^{\mathrm{E}}$ обращает в нуль правую часть \eqref{eq:Lambda-es-xi}, если само равенство выполнено. Явные интегралы главы~\ref{sec:es-beta} такое обнуление не доказывают: горизонтальный остаток в них не вошёл. Натуральную часть поправляют секступоли и нелинейность дефлектора. Вклад дефлектора фиксирован показателем $n$.

В магнитном выводе интеграл одних наклонов $\oint\gamma\,\mathrm{d}s$ при $\xi=0$ в нуль не обращается: его компенсирует сдвиг центра. Секступоль в линейное уравнение для $D_1$ не входит. Сила секступоля квадратична по смещению, и на опорной орбите, проходящей через его центр, линейного члена нет. На орбите $x=D_1\delta$ та же сила пропорциональна $\delta^2$ и стоит в правой части уравнения для $D_2$. Поэтому $\alpha_1$ и $\eta_1$ от этих секступолей зависят. Величина сдвига $\eta_1$ в главах~\ref{sec:mag-base}--\ref{sec:es-beta} не получена. Коэффициенты $\eta_0$ и $\alpha_0$ коррекция не меняет.

После коррекции в \eqref{eq:deq-L} и \eqref{eq:deq-es-L} остаётся продольный член
\begin{equation}
\label{eq:cmp-deq-xi0}
\delta_\eq=-\frac{\eta_1\delta_m^2}{2\eta_0}.
\end{equation}
Во втором порядке по фазе при члене с $\eta_1$ стоит множитель $(1-\hat\varphi^2/32)$ из \eqref{eq:deq-ph-L}. На $\Lambda$ этот множитель не действует, поэтому при $\Lambda=0$ он остаётся у члена с $\eta_1$.

Зависимость $\Lambda$ от $\delta$ в продольных уравнениях отброшена как член порядка $\delta J$. Условие $\xi=0$ относится к той же точке $\delta=0$, в которой $\Lambda$ вычислено. Производная тюна второго порядка по $\delta$ в сохранённый $\Lambda$ не входит.

\subsection{Сравнение коэффициентов}

Магнитная форма --- выражение \eqref{eq:tau}, в которое подставлены геометрические $\alpha_0$ и $\alpha_1$ рассматриваемого кольца. Для электростатического кольца это его собственные $\alpha_0^{\mathrm{E}}$ и $\alpha_1^{\mathrm{E}}$.

Линейные коэффициенты
\begin{equation}
\label{eq:cmp-eta0}
\eta_0^{\mathrm{M}}
=\alpha_0^{\mathrm{M}}-\frac{1}{\gamma_s^2},
\qquad
\eta_0^{\mathrm{E}}
=\Bigl(1+\frac{1}{\gamma_s^2}\Bigr)\alpha_0^{\mathrm{E}}-\frac{1}{\gamma_s^2}.
\end{equation}
Первое есть \eqref{eq:tau}, второе --- \eqref{eq:eta0-es}. Разность
\begin{equation}
\label{eq:cmp-eta0-diff}
\eta_0^{\mathrm{E}}-\eta_0^{\mathrm{M}}
=\bigl(\alpha_0^{\mathrm{E}}-\alpha_0^{\mathrm{M}}\bigr)
+\frac{\alpha_0^{\mathrm{E}}}{\gamma_s^2}.
\end{equation}
Первое слагаемое --- разность удлинений орбиты. Функции $D_1^{\mathrm{E}}$ и $D_1^{\mathrm{M}}$ удовлетворяют \eqref{eq:cmp-D-es} и \eqref{eq:cmp-D-mag}. Второе слагаемое пропорционально $1/\gamma_s^2$. Его источник --- скорость в потенциале. На внешней орбите $\psi>0$, кинетическая энергия ниже, частица медленнее. В ряде \eqref{eq:vsv-es} это член $(D_1/\rho)\delta/\gamma_s^2$, и его среднее есть $\alpha_0^{\mathrm{E}}/\gamma_s^2$.

Квадратичный магнитный коэффициент есть \eqref{eq:tau},
\begin{equation}
\label{eq:cmp-eta1-mag}
\eta_1^{\mathrm{M}}
=\alpha_1^{\mathrm{M}}-\frac{\alpha_0^{\mathrm{M}}}{\gamma_s^2}
+\frac{3\gamma_s^2-1}{2\gamma_s^4}.
\end{equation}
Электростатический \eqref{eq:eta1-es} на тех же обозначениях \eqref{eq:cmp-alpha} равен магнитной форме на $\alpha^{\mathrm{E}}$ плюс поправка:
\begin{equation}
\label{eq:cmp-eta1}
\begin{aligned}
\eta_1^{\mathrm{E}}
&=\alpha_1^{\mathrm{E}}-\frac{\alpha_0^{\mathrm{E}}}{\gamma_s^2}
+\frac{3\gamma_s^2-1}{2\gamma_s^4}\\
&\quad+\frac{1}{\gamma_s^2}\Biggl[
\Biggl\langle\frac{D_2}{\rho}\Biggr\rangle_C
+\frac{5-n}{2}
\Biggl\langle\Bigl(\frac{D_1}{\rho}\Bigr)^2\Biggr\rangle_C
-3\alpha_0^{\mathrm{E}}
\Biggr].
\end{aligned}
\end{equation}
В скобке стоят $D_1$ и $D_2$ электростатического кольца. Наклон $\langle D_1'^2\rangle_C$ в поправку не входит: он уже учтён в $\alpha_1^{\mathrm{E}}$, а произведение наклона с поправкой скорости есть третий порядок, как в выводе \eqref{eq:eta1-es}.

Член $\langle D_2/\rho\rangle_C$ в скобке --- линейный потенциал на смещении $D_2\delta^2$, тот же механизм, что $\alpha_0^{\mathrm{E}}/\gamma_s^2$ в \eqref{eq:cmp-eta0-diff}. Геометрическая часть $D_2$ уже сидит в $\alpha_1^{\mathrm{E}}$. Член $(5-n)/2\,\langle(D_1/\rho)^2\rangle_C$ собран в главе~\ref{sec:es-base} из трёх вкладов: $-n/2$ от квадратичного потенциала, $+1$ от произведения удлинения дуги и линейного замедления, $+3/2$ от квадрата $3(\delta-\psi)^2/(2\gamma_s^2)$. Слагаемое $-3\alpha_0^{\mathrm{E}}$ внутри скобки, то есть $-3\alpha_0^{\mathrm{E}}/\gamma_s^2$ в $\eta_1^{\mathrm{E}}$, --- перекрёстный член того же квадрата. На орбите $\delta-\psi=(1-D_1/\rho)\delta$, и $(1-D_1/\rho)^2$ содержит $-2D_1/\rho$. В магнитном кольце $\psi=0$ вдоль всей орбиты, поэтому такого перекрёстного члена нет. Общий кусок $(3\gamma_s^2-1)/(2\gamma_s^4)$ в разности коэффициентов сокращается.

Разность коэффициентов двух колец поэтому содержит и разные $\alpha$, и потенциал:
\begin{equation}
\label{eq:cmp-eta1-diff}
\begin{aligned}
\eta_1^{\mathrm{E}}-\eta_1^{\mathrm{M}}
&=\bigl(\alpha_1^{\mathrm{E}}-\alpha_1^{\mathrm{M}}\bigr)
-\frac{\alpha_0^{\mathrm{E}}-\alpha_0^{\mathrm{M}}}{\gamma_s^2}\\
&\quad+\frac{1}{\gamma_s^2}\Biggl[
\Biggl\langle\frac{D_2}{\rho}\Biggr\rangle_C
+\frac{5-n}{2}
\Biggl\langle\Bigl(\frac{D_1}{\rho}\Bigr)^2\Biggr\rangle_C
-3\alpha_0^{\mathrm{E}}
\Biggr].
\end{aligned}
\end{equation}
Первая строка --- разность геометрических $\alpha_1$. Вторая и третья строки пропорциональны $1/\gamma_s^2$. Вторая --- разность членов $-\alpha_0/\gamma_s^2$ из \eqref{eq:tau} на двух орбитах. Третья --- потенциал электростатического кольца, та же скобка, что в \eqref{eq:cmp-eta1}. В $\alpha_1^{\mathrm{E}}$ входит дисперсия электростатического кольца, поэтому зависимость $D_1$ и $D_2$ от $1/\gamma_s^2$ входит и в первую строку.

Бетатронное запаздывание сравнивается так же. Геометрическая часть --- наклоны и сдвиг орбиты. Это выражение \eqref{eq:Lambda-slopes} на функциях Твисса и на $x_{\mathrm{nl}}$ электростатического кольца. По \eqref{eq:xp2} оно равно
\begin{equation}
\label{eq:cmp-geom}
\begin{aligned}
&\frac{1}{C_s}\oint\Biggl(
\frac{x_{\mathrm{nl}}}{\rho}
+\frac{\langle x_\beta'^2\rangle_\psi+\langle y_\beta'^2\rangle_\psi}{2}
\Biggr)\mathrm{d}s
\\
&\qquad=
\frac{1}{C_s}\Biggl[
\frac{J_x}{2}\oint\gamma_x\,\mathrm{d}s
+\frac{J_y}{2}\oint\gamma_y\,\mathrm{d}s
+\oint\frac{x_{\mathrm{nl}}}{\rho}\,\mathrm{d}s
\Biggr].
\end{aligned}
\end{equation}
Магнитное $\Lambda^{\mathrm{M}}$ есть та же геометрическая часть на оптике магнитного кольца. Вычитание правой части \eqref{eq:cmp-geom} из \eqref{eq:Lambda-es-J} оставляет часть \eqref{eq:Lambda-es}, пропорциональную $1/\gamma_s^2$:
\begin{equation}
\label{eq:cmp-Lambda}
\begin{aligned}
&\Lambda^{\mathrm{E}}
-\frac{1}{C_s}\Biggl[
\frac{J_x}{2}\oint\gamma_x\,\mathrm{d}s
+\frac{J_y}{2}\oint\gamma_y\,\mathrm{d}s
+\oint\frac{x_{\mathrm{nl}}}{\rho}\,\mathrm{d}s
\Biggr]
\\
&\quad=\frac{1}{C_s\gamma_s^2}
\Biggl(
\frac{J_x}{2}\oint\frac{(5-n)\beta_x}{\rho^2}\,\mathrm{d}s
+\frac{J_y}{2}\oint\frac{(n-1)\beta_y}{\rho^2}\,\mathrm{d}s
+\oint\frac{x_{\mathrm{nl}}}{\rho}\,\mathrm{d}s
\Biggr).
\end{aligned}
\end{equation}
Первые два интеграла --- те же вклады потенциала, что в \eqref{eq:Lambda-es}, теперь на средних $\langle x_\beta^2\rangle_\psi=J_x\beta_x$ и $\langle y_\beta^2\rangle_\psi=J_y\beta_y$. Коэффициенты при этих средних --- $(5-n)/(2\gamma_s^2)$ и $(n-1)/(2\gamma_s^2)$. Вертикальное отклонение входит в свой член из \eqref{eq:psi}. Третий интеграл --- скорость на сдвиге центра. Удлинение пути от $x_{\mathrm{nl}}$ уже стоит в \eqref{eq:cmp-geom}, а множитель $1/\gamma_s^2$ берётся из \eqref{eq:v-delta0}.

Формула \eqref{eq:cmp-Lambda} --- потенциал на оптике электростатического кольца. В разности $\Lambda^{\mathrm{E}}-\Lambda^{\mathrm{M}}$ к этому потенциалу добавляются разные $\beta_x$, $\gamma_x$ и $x_{\mathrm{nl}}$ двух колец.

\subsection{Эффективный лоренц-фактор}

В главах~\ref{sec:mag-base}, \ref{sec:mag-beta}, \ref{sec:es-base} и~\ref{sec:es-beta} величина $\gamma_{\mathrm{eff}}$ --- среднее лоренц-фактора за синхротронный период там, где потенциал равен нулю. Усреднение ряда \eqref{eq:gamma-delta} оставляет $\gamma_s$, член $\beta_s^2\gamma_s\delta_\eq$ и член $\beta_s^2\langle\delta^2\rangle/(2\gamma_s)$. Множитель $\beta_s^2$ переводит $\delta$ в приращение $\gamma$, остаток есть $O(\delta_m^4)$. В первом порядке по фазе $\langle\delta^2\rangle=\delta_m^2/2$, во втором --- $(\delta_m^2/2)(1-\hat\varphi^2/32)$. Член с $\Lambda$ входит в это среднее только через $\delta_\eq$. Та же запись собрана в \eqref{eq:cmp-master}.

Формула (10) коллег сохраняет свободный член и член с $\delta_\eq$:
\begin{equation}
\label{eq:cmp-colleagues}
\gamma_{\mathrm{eff}}
=\gamma_s+\beta_s^2\gamma_s\,\delta_\eq.
\end{equation}
Разность полного среднего и \eqref{eq:cmp-colleagues} равна $\beta_s^2\langle\delta^2\rangle/(2\gamma_s)$. Оба сдвига второго порядка по $\delta_m$, поскольку само $\delta_\eq$ порядка $\delta_m^2$. Отношение отброшенного члена к модулю сохранённого есть
\begin{equation}
\label{eq:cmp-ratio}
\frac{\langle\delta^2\rangle}{2\gamma_s^2\,|\delta_\eq|}
=\frac{|\eta_0|\,\delta_m^2}
{4\gamma_s^2\bigl|\eta_1\delta_m^2/2+\Lambda\bigr|}.
\end{equation}
Равенство \eqref{eq:cmp-ratio} записано при $\sin\varphi\approx\varphi$. При $\Lambda=0$ правая часть есть $|\eta_0|/(2|\eta_1|\gamma_s^2)$, как после \eqref{eq:gamma-eff}. Множитель $(1-\hat\varphi^2/32)$ в отношении сокращается, и амплитуда $\delta_m$ в него не входит. Отношение порядка единицы, когда $|\eta_1|$ порядка $|\eta_0|/\gamma_s^2$, и мало, когда $|\eta_1|$ много больше этой величины. Множитель $1/\gamma_s^2$ в \eqref{eq:cmp-ratio} происходит из кинематического члена $\beta_s^2\langle\delta^2\rangle/(2\gamma_s)$, отнесённого к сдвигу $\beta_s^2\gamma_s\delta_\eq$. При $\eta_1=0$ в знаменателе стоит $|\Lambda|$. Условие $|\Lambda/(\eta_0\delta_m)|\ll 1$ само по себе это отношение малым не делает. После коррекции хроматичности сравнение снова идёт через $\eta_1$, уже сдвинутое секступолями. Величина сдвига в записке не вычислена.

В электростатическом кольце среднее \eqref{eq:cmp-master} относится к участкам с $\Phi=0$. Внутри дефлектора в фиксированной точке $s$ среднее по бетатронной фазе и по синхротронному периоду есть \eqref{eq:gamma-loc-avg}. Из $\gamma_{\mathrm{eff}}$ вычитается $\beta_s^2\gamma_s\langle\psi\rangle_s$, а $\langle\psi\rangle_s$ берётся из \eqref{eq:psi-avg-s}. Первое слагаемое \eqref{eq:psi-avg-s} даёт в вычитаемом член $\beta_s^2\gamma_s(D_1/\rho)\delta_\eq$. Множитель $D_1/\rho$ берётся из орбиты в этой точке. Внутри дефлектора $D_1$ задано уравнением \eqref{eq:cmp-D-es}. Одно число $\langle\psi\rangle$ по всему дефлектору из \eqref{eq:psi-avg-s} не получается: для него нужен вес вдоль $s$, а уравнение продольного движения этот вес не фиксирует.

Записка вычисляет эти средние и не выводит частоту прецессии спина. Формула \eqref{eq:cmp-colleagues} совпадает с частью $\langle\gamma\rangle$, линейной по $\delta_\eq$, вне дефлекторов и в магнитном кольце. Внутри дефлектора усредняется локальный кинетический лоренц-фактор \eqref{eq:gamma-loc-avg}.

Сводка собрана в одной записи. Средний квадрат берётся на решении соответствующего порядка: $\langle\delta^2\rangle=\delta_m^2/2$ при $\sin\varphi\approx\varphi$ и $\langle\delta^2\rangle=(\delta_m^2/2)(1-\hat\varphi^2/32)$ при учёте $\varphi^3$. Множитель $(1-\hat\varphi^2/32)$ не умножает $\Lambda$.
\begin{equation}
\label{eq:cmp-master}
\begin{aligned}
\delta_\eq
&=-\frac{1}{\eta_0}\bigl(\eta_1\langle\delta^2\rangle+\Lambda\bigr),\\[2pt]
\gamma_{\mathrm{eff}}
&=\gamma_s+\beta_s^2\gamma_s\,\delta_\eq
+\frac{\beta_s^2}{2\gamma_s}\langle\delta^2\rangle.
\end{aligned}
\end{equation}
Остаток в $\gamma_{\mathrm{eff}}$ имеет порядок $\delta_m^4$. Какой коэффициент подставить, указано в таблице. Для электростатического кольца при $\Lambda^{\mathrm{E}}=0$ первая строка есть \eqref{eq:deq-ph} с $\eta^{\mathrm{E}}$, как в подразделе второго порядка по фазе главы~\ref{sec:es-base}. Строки $\gamma_{\mathrm{eff}}$ получаются усреднением \eqref{eq:gamma-delta} с тем же $\langle\delta^2\rangle$. При ненулевом $\Lambda^{\mathrm{E}}$ строки $\delta_\eq$ и $\gamma_{\mathrm{eff}}$ имеют вид \eqref{eq:deq-ph-L} с коэффициентами этого кольца. Глава~\ref{sec:es-beta} выводит \eqref{eq:deq-es-L} при $\sin\varphi\approx\varphi$.

\begin{center}
\begingroup
\newcommand{\cmpcell}[1]{\parbox[t]{0.36\textwidth}{\raggedright #1\strut}}
\renewcommand{\arraystretch}{1.25}
\begin{tabular}{@{}llp{0.28\textwidth}@{}}
\hline
случай & $\eta$, $\Lambda$ & формулы \\
\hline
\cmpcell{магнитное, без $\Lambda$, $\sin\varphi\approx\varphi$}
& $\eta^{\mathrm{M}}$, $\Lambda=0$
& \eqref{eq:deq}, \eqref{eq:gamma-eff} \\
\cmpcell{магнитное, без $\Lambda$, до $\varphi^3$}
& $\eta^{\mathrm{M}}$, $\Lambda=0$
& \eqref{eq:deq-ph}, \eqref{eq:deq-ph-L} \\
\cmpcell{магнитное, с $\Lambda$, $\sin\varphi\approx\varphi$}
& $\eta^{\mathrm{M}}$, $\Lambda^{\mathrm{M}}$
& \eqref{eq:deq-L}, \eqref{eq:gamma-L} \\
\cmpcell{магнитное, с $\Lambda$, до $\varphi^3$}
& $\eta^{\mathrm{M}}$, $\Lambda^{\mathrm{M}}$
& \eqref{eq:deq-ph-L} \\
\cmpcell{электростатическое, без $\Lambda$, $\sin\varphi\approx\varphi$}
& $\eta^{\mathrm{E}}$, $\Lambda=0$
& \eqref{eq:deq-es}, \eqref{eq:gamma-eff-es} \\
\cmpcell{электростатическое, без $\Lambda$, до $\varphi^3$}
& $\eta^{\mathrm{E}}$, $\Lambda=0$
& \eqref{eq:deq-ph}, \eqref{eq:deq-ph-L} \\
\cmpcell{электростатическое, с $\Lambda$, $\sin\varphi\approx\varphi$}
& $\eta^{\mathrm{E}}$, $\Lambda^{\mathrm{E}}$
& \eqref{eq:deq-es-L}, \eqref{eq:gamma-es-L} \\
\cmpcell{электростатическое, с $\Lambda$, до $\varphi^3$}
& $\eta^{\mathrm{E}}$, $\Lambda^{\mathrm{E}}$
& \eqref{eq:deq-ph-L} \\
\hline
\end{tabular}
\endgroup
\end{center}
Ссылка на \eqref{eq:deq-ph-L} означает все три строки формулы с коэффициентами данного кольца; при $\Lambda=0$ эта постоянная равна нулю. Внутри дефлектора вместо $\gamma_{\mathrm{eff}}$ стоит \eqref{eq:gamma-loc-avg}.


\clearpage
\appendix
\section{Метод Боголюбова--Митропольского}
\label{app:bm}

\subsection{Уравнение с квадратичной нелинейностью}
\label{sec:bm-quad}

Уравнение \eqref{eq:ode} с произвольной постоянной $c_0$ того же порядка, что и квадратичные члены, записывается так:
\begin{equation}
\label{eq:bm}
\ddot\delta+\omega^2\delta
=\varepsilon\Biggl(
c_0
+c_2\delta^2
+c_d(\dot\delta)^2
\Biggr),
\qquad
\omega=\omega_s,
\end{equation}
\begin{equation}
\label{eq:bm-coeff}
c_2=-\frac{\omega^2}{\eta_0}\Bigl(\eta_1-\frac{\eta_0}{\gamma_s^2}\Bigr),
\qquad
c_d=-\frac{1}{\gamma_s^2}.
\end{equation}
В главе~\ref{sec:mag-base} выполнено $c_0=0$. Если к $\tau$ добавляется постоянная $\Lambda$ одного порядка с $\delta^2$, в скобке \eqref{eq:ddelta-raw} появляется слагаемое $\Lambda$. Здесь $\kappa=K\cos\varphi_s$, и по \eqref{eq:ws} выполнено $\kappa\omega_\rf=-\omega^2/\eta_0$. Умножение на этот множитель даёт на правой части \eqref{eq:ode} постоянную $-\omega^2\Lambda/\eta_0$. Перекрёстный член $\Lambda\delta$ имеет третий порядок по амплитуде и отбрасывается, откуда
\begin{equation}
\label{eq:c0-L}
c_0=-\frac{\omega^2\Lambda}{\eta_0}.
\end{equation}
Параметр $\varepsilon$ отмечает порядок малости и в ответе полагается равным единице~\cite{Bogoliubov1961}.

Решение ищется в виде
\begin{equation}
\label{eq:ansatz}
\begin{aligned}
\delta&=a\cos\psi+\varepsilon u_1(a,\psi)+\varepsilon^2 u_2(a,\psi),\\
\dot a&=\varepsilon A_1+\varepsilon^2 A_2,
\qquad
\dot\psi&=\omega+\varepsilon B_1+\varepsilon^2 B_2.
\end{aligned}
\end{equation}
Функции $u_k$ периодичны по $\psi$ с периодом $2\pi$ и не содержат первой гармоники: коэффициенты при $\cos\psi$ и $\sin\psi$ равны нулю. Этим фиксирован смысл $a$ и $\psi$: вся первая гармоника собрана в $a\cos\psi$.

Подстановка в левую часть до порядка $\varepsilon$ считается явно. Ведущий член $a\cos\psi$ имеет производную
\begin{equation}
\label{eq:v0}
\frac{\mathrm{d}}{\mathrm{d}t}(a\cos\psi)
=-a\omega\sin\psi
+\varepsilon(A_1\cos\psi-a B_1\sin\psi)+O(\varepsilon^2).
\end{equation}
Здесь использовано $\dot a=\varepsilon A_1+O(\varepsilon^2)$ и $\dot\psi=\omega+\varepsilon B_1+O(\varepsilon^2)$. Производная поправки $\varepsilon u_1$ на том же порядке равна $\varepsilon\omega\,\partial_\psi u_1$, потому что $\dot a$ уже содержит $\varepsilon$, а $\dot\psi$ заменяется на $\omega$: следующий порядок есть $O(\varepsilon^2)$. Сумма
\begin{equation}
\label{eq:v1def}
\dot\delta
=-a\omega\sin\psi
+\varepsilon v_1+O(\varepsilon^2),
\qquad
v_1=A_1\cos\psi-a B_1\sin\psi+\omega\,\partial_\psi u_1.
\end{equation}
Вторая производная ведущего куска $-a\omega\sin\psi$ равна
\begin{equation}
\label{eq:d2v0}
\frac{\mathrm{d}}{\mathrm{d}t}(-a\omega\sin\psi)
=-a\omega^2\cos\psi
+\varepsilon(-A_1\omega\sin\psi-a\omega B_1\cos\psi)+O(\varepsilon^2).
\end{equation}
Вторая производная куска $\varepsilon v_1$ равна $\varepsilon\omega\,\partial_\psi v_1+O(\varepsilon^2)$. Раскрытие
\begin{equation}
\label{eq:dv1}
\partial_\psi v_1=-A_1\sin\psi-a B_1\cos\psi+\omega\,\partial_{\psi\psi}u_1
\end{equation}
даёт ещё один такой же вклад $-A_1\omega\sin\psi-a\omega B_1\cos\psi$. Два одинаковых вклада складываются, и коэффициент равен двум. Вместе с $\omega^2\cdot\varepsilon u_1$ член порядка $\varepsilon$ в $\ddot\delta+\omega^2\delta$ есть
\begin{equation}
\label{eq:lhs1}
\omega^2(\partial_{\psi\psi}u_1+u_1)-2\omega A_1\sin\psi-2a\omega B_1\cos\psi.
\end{equation}
На каждом следующем порядке $\varepsilon^k$ новые неизвестные $A_k$, $B_k$, $u_k$ входят в левую часть так же: $\dot a$ и $\dot\psi$ стоят в производной $a\cos\psi$ линейно, а $u_k$ входит в $\delta$ и во вторую производную линейно. Произведения уже найденных младших поправок переносятся в правую часть. Поэтому на порядке $\varepsilon^k$
\begin{equation}
\label{eq:ordk}
\omega^2(\partial_{\psi\psi}u_k+u_k)-2\omega A_k\sin\psi-2a\omega B_k\cos\psi=f_k,
\end{equation}
где $f_k$ --- коэффициент при $\varepsilon^k$ после этого переноса.

Правило гармоник следует из прямой подстановки: $\partial_{\psi\psi}\cos(m\psi)=-m^2\cos(m\psi)$, значит
\begin{equation}
\label{eq:harm-rule}
(\partial_{\psi\psi}+1)\cos(m\psi)=(1-m^2)\cos(m\psi).
\end{equation}
При $m\neq 1$ оператор обратим. Если правая часть уравнения $\omega^2(\partial_{\psi\psi}+1)u=F_m\cos(m\psi)$ содержит такую гармонику, коэффициент в $u$ равен $F_m/[\omega^2(1-m^2)]$. При $m=1$ выполнено $1-m^2=0$, оператор аннулирует $\cos\psi$ и $\sin\psi$, и эти члены правой части определяют $A_k$ и $B_k$, а в $u_k$ не входят.

На первом порядке правая часть \eqref{eq:bm} берётся на ведущем решении $\delta=a\cos\psi$, $\dot\delta=-a\omega\sin\psi$. Поправки порядка $\varepsilon$ внутри $\delta^2$ и $(\dot\delta)^2$ попадают в $\varepsilon^2$. Тождества $\cos^2\psi=(1+\cos 2\psi)/2$ и $\sin^2\psi=(1-\cos 2\psi)/2$ дают
\begin{equation}
\label{eq:f1}
\begin{aligned}
f_1
&=c_0+c_2 a^2\frac{1+\cos 2\psi}{2}
+c_d a^2\omega^2\frac{1-\cos 2\psi}{2}\\
&=c_0+\frac{a^2}{2}(c_2+c_d\omega^2)
+\frac{a^2}{2}(c_2-c_d\omega^2)\cos 2\psi.
\end{aligned}
\end{equation}
Гармоник $\cos\psi$ и $\sin\psi$ в \eqref{eq:f1} нет: их коэффициенты равны нулю. Левая часть \eqref{eq:lhs1} совпадает с $f_1$, поэтому коэффициенты при $\sin\psi$ и $\cos\psi$ слева тоже равны нулю. В \eqref{eq:lhs1} коэффициент при $\sin\psi$ равен $-2\omega A_1$, коэффициент при $\cos\psi$ равен $-2a\omega B_1$. При $\omega\neq 0$ и $a\neq 0$ оба коэффициента обращаются в нуль только при $A_1=0$ и $B_1=0$. Амплитуда на этом порядке постоянна, сдвига частоты нет.

Оставшиеся гармоники определяют $u_1=U_0+U_2\cos 2\psi$. Вторая производная постоянной равна нулю, так что $(\partial_{\psi\psi}+1)U_0=U_0$ и
\begin{equation}
\label{eq:U0}
\omega^2 U_0=c_0+\frac{a^2}{2}(c_2+c_d\omega^2).
\end{equation}
На гармонике $\cos 2\psi$ номер $m=2$. По правилу \eqref{eq:harm-rule} множитель $1-m^2$ равен $1-4=-3$, так что $\omega^2(\partial_{\psi\psi}+1)$ умножает коэффициент $U_2$ на $-3\omega^2$. Это произведение равно коэффициенту при $\cos 2\psi$ в \eqref{eq:f1}, поэтому
\begin{equation}
\label{eq:U2}
-3\omega^2 U_2=\frac{a^2}{2}(c_2-c_d\omega^2),
\qquad
U_2=-\frac{a^2(c_2-c_d\omega^2)}{6\omega^2}.
\end{equation}

Коэффициенты \eqref{eq:bm-coeff} подставляются явно. Сумма
\begin{equation}
\label{eq:c-plus}
c_2+c_d\omega^2
=-\frac{\omega^2}{\eta_0}\eta_1+\frac{\omega^2}{\gamma_s^2}-\frac{\omega^2}{\gamma_s^2}
=-\frac{\omega^2\eta_1}{\eta_0}.
\end{equation}
Два слагаемых $\omega^2/\gamma_s^2$ взаимно уничтожаются: первое возникает при раскрытии $(\omega^2/\eta_0)\cdot(\eta_0/\gamma_s^2)$ внутри $c_2$, второе равно $c_d\omega^2$. Деление \eqref{eq:U0} на $\omega^2$ даёт
\begin{equation}
\label{eq:U0-simp}
U_0=\frac{c_0}{\omega^2}-\frac{\eta_1 a^2}{2\eta_0}.
\end{equation}
Разность
\begin{equation}
\label{eq:c-minus}
c_2-c_d\omega^2
=-\frac{\omega^2\eta_1}{\eta_0}+\frac{2\omega^2}{\gamma_s^2},
\end{equation}
потому что $-c_d\omega^2=+\omega^2/\gamma_s^2$ добавляется к такому же члену из $(\omega^2/\eta_0)\cdot(\eta_0/\gamma_s^2)$. Подстановка в \eqref{eq:U2} даёт
\begin{equation}
\label{eq:U2-simp}
U_2=\frac{a^2}{6}\Bigl(\frac{\eta_1}{\eta_0}-\frac{2}{\gamma_s^2}\Bigr).
\end{equation}
При $\varepsilon=1$, с $A_1=B_1=0$ и с началом отсчёта, убирающим постоянную в $\psi$,
\begin{equation}
\label{eq:bm-sol}
\begin{aligned}
\delta
&=\frac{c_0}{\omega_s^2}-\frac{\eta_1 a^2}{2\eta_0}
+a\cos\psi
+\frac{a^2}{6}\Bigl(\frac{\eta_1}{\eta_0}-\frac{2}{\gamma_s^2}\Bigr)\cos 2\psi,\\
\psi&=\omega_s t.
\end{aligned}
\end{equation}
Сдвиг частоты появляется на следующем порядке метода из проекции на $\cos\psi$. Поправка этого порядка, пропорциональная $a^2$, в физической задаче требует члена $\delta^3$ в $\tau$ и здесь не выписывается.

\subsection{Уравнение с кубической нелинейностью}
\label{sec:bm-cubic}

Уравнение \eqref{eq:phi-ode} записывается так:
\begin{equation}
\label{eq:bm-c}
\ddot\varphi+\omega^2\varphi=\varepsilon c_3\varphi^3,
\qquad
\omega=\omega_s,
\qquad
c_3=\frac{\omega^2}{6}.
\end{equation}
Перенос члена $\omega^2\varphi^3/6$ из левой части \eqref{eq:phi-ode} в правую даёт $c_3\varphi^3$ при $\varepsilon=1$. Параметр $\varepsilon$ отмечает порядок малости и в ответе полагается равным единице, как в \S\ref{sec:bm-quad}.

Схема \eqref{eq:ansatz}, \eqref{eq:ordk} применяется к $\varphi$ вместо $\delta$. В этой подстановке $a$ есть коэффициент при $\cos\psi$ в фазе. Функция $u_1$ по условию схемы не содержит $\cos\psi$ и $\sin\psi$.

На первом порядке правая часть берётся на $\varphi=a\cos\psi$. Тождество $\cos 3\psi=4\cos^3\psi-3\cos\psi$ решается относительно куба:
\begin{equation}
\label{eq:cos3id}
\cos^3\psi=\frac{3\cos\psi+\cos 3\psi}{4}.
\end{equation}
Подстановка \eqref{eq:cos3id} даёт
\begin{equation}
\label{eq:f1c}
f_1
=c_3 a^3\cos^3\psi
=\frac{3c_3 a^3}{4}\cos\psi
+\frac{c_3 a^3}{4}\cos 3\psi.
\end{equation}
Коэффициент при $\sin\psi$ в \eqref{eq:f1c} равен нулю: разложение \eqref{eq:cos3id} его не содержит. В левой части \eqref{eq:ordk} коэффициент при $\sin\psi$ равен $-2\omega A_1$. При $\omega\neq 0$ равенство коэффициентов даёт $A_1=0$.

Коэффициент при $\cos\psi$ в \eqref{eq:f1c} равен $3c_3 a^3/4$. В левой части \eqref{eq:ordk} он равен $-2a\omega B_1$: по \eqref{eq:harm-rule} при $m=1$ оператор $\partial_{\psi\psi}+1$ аннулирует $\cos\psi$, а первая гармоника в $u_1$ запрещена анзацем \eqref{eq:ansatz}. При $a\neq 0$ и $\omega\neq 0$
\begin{equation}
\label{eq:B1c}
-2a\omega B_1=\frac{3c_3 a^3}{4},
\qquad
B_1=-\frac{3c_3 a^2}{8\omega}.
\end{equation}

На гармонике $\cos 3\psi$ номер $m=3$. Правило \eqref{eq:harm-rule} даёт $1-m^2=1-9=-8$, так что $\omega^2(\partial_{\psi\psi}+1)$ умножает коэффициент $U_3$ на $-8\omega^2$. Это произведение равно коэффициенту $c_3 a^3/4$ в \eqref{eq:f1c}, поэтому
\begin{equation}
\label{eq:U3c}
-8\omega^2 U_3=\frac{c_3 a^3}{4},
\qquad
U_3=-\frac{c_3 a^3}{32\omega^2},
\qquad
u_1=U_3\cos 3\psi.
\end{equation}
Подстановка $c_3=\omega^2/6$ сокращает $\omega^2$. В $B_1$ получается $3(\omega^2/6)a^2/(8\omega)=\omega a^2/16$, знак минус сохраняется. В $U_3$ получается $(\omega^2/6)a^3/(32\omega^2)=a^3/192$, знак минус сохраняется:
\begin{equation}
\label{eq:B1U3}
B_1=-\frac{\omega a^2}{16},
\qquad
U_3=-\frac{a^3}{192}.
\end{equation}
При $\varepsilon=1$ частота есть $\dot\psi=\omega+B_1$, и
\begin{equation}
\label{eq:phi-a}
\varphi
=a\cos\psi-\frac{a^3}{192}\cos 3\psi,
\qquad
\Omega=\omega\Bigl(1-\frac{a^2}{16}\Bigr).
\end{equation}
Здесь $a$ --- коэффициент первой гармоники, не полуразмах.

В главе~\ref{sec:mag-phase} амплитуда задана полуразмахом; следующие равенства только переводят $a$ в $\hat\varphi$. При $\psi=0$ выполнено $\cos\psi=\cos 3\psi=1$, поэтому значение \eqref{eq:phi-a} равно $a-a^3/192$. При $\psi=\pi$ оно противоположно. Производная по $\psi$ равна $-a\sin\psi+(a^3/64)\sin 3\psi$: множитель $3/192$ равен $1/64$. Вынос $\sin\psi$ оставляет множитель $1-(a^2/64)(\sin 3\psi/\sin\psi)$. Отношение $\sin 3\psi/\sin\psi=3-4\sin^2\psi$ ограничено, и при $a^2\ll 1$ этот множитель в нуль не обращается. Нули производной только при $\sin\psi=0$. Полуразмах равен значению при $\psi=0$:
\begin{equation}
\label{eq:phihat-a}
\hat\varphi=a\Bigl(1-\frac{a^2}{192}\Bigr).
\end{equation}
Обращение с отбрасыванием $O(a^5)$ есть $a=\hat\varphi+\hat\varphi^3/192$. Подстановка этой формулы обратно в $a-a^3/192$ возвращает $\hat\varphi$, потому что $a^3=\hat\varphi^3+O(\hat\varphi^5)$. Подстановка того же $a$ в \eqref{eq:phi-a} даёт
\begin{equation}
\label{eq:phi-hat-cos}
\varphi
=\hat\varphi\cos\psi
+\frac{\hat\varphi^3}{192}\bigl(\cos\psi-\cos 3\psi\bigr),
\qquad
\frac{\Omega}{\omega_s}=1-\frac{\hat\varphi^2}{16}.
\end{equation}
В частоте $a^2=\hat\varphi^2+2\hat\varphi\cdot(\hat\varphi^3/192)+O(\hat\varphi^6)=\hat\varphi^2+O(\hat\varphi^4)$, поэтому замена $a^2$ на $\hat\varphi^2$ не меняет член $\hat\varphi^2/16$. Остаток в $\varphi$ есть $O(\hat\varphi^5)$.



\clearpage
\bibliographystyle{unsrt}
\bibliography{references}

\end{document}
